\color{unidarkgreen}{\section[(обязательное) Уставки функций РЗА]{(обязательное)\\Уставки функций РЗА}\label{app:set}}
\color{black}

\setlength{\extrarowheight}{0.06cm}



\par
Уставки ДТЗ представлены в таблице \ref{app_set:dzt}.

\startTable{Параметры для настройки ДТЗ}{app_set:dzt}

%===m>ДТЗ|general+-
\multicolumn{5}{|c|}{ ДТЗ } \\ \hline 
\raggedright ДТЗ фВ & \centering Выведено \\ Введено & \centering --- & \centering --- & \centering \arraybackslash Выведено \\
\hline
\raggedright Iср диф & \centering 0,100 ... 2,000 & \centering о.е. & \centering 0,001 & \centering \arraybackslash 0,200 \\
\hline
\raggedright Кв(Iср диф) & \centering 0,00 ... 1,00 & \centering о.е. & \centering 0,01 & \centering \arraybackslash 0,80 \\
\hline
\multicolumn{5}{|c|}{ КЦТ } \\ \hline 
\raggedright Iнеиспр ДТЗ & \centering 0,04 ... 30,00 & \centering о.е. & \centering 0,01 & \centering \arraybackslash 0,10 \\
\hline
\raggedright Кв (Iнеиспр) & \centering 0,00 ... 1,00 & \centering  & \centering 0,01 & \centering \arraybackslash 0,95 \\
\hline
\raggedright Тнеиспр ДТЗ & \centering 0 ... 300000 & \centering мс & \centering 5 & \centering \arraybackslash 10000 \\
\hline
\raggedright Контр изм цеп ДТЗ & \centering Введено \\ Выведено & \centering --- & \centering --- & \centering \arraybackslash Введено \\
\hline
\raggedright фВ & \centering Введено \\ Выведено & \centering --- & \centering --- & \centering \arraybackslash Введено \\
\hline
\multicolumn{5}{|c|}{ ДТЗт } \\ \hline 
\raggedright Iср диф & \centering 0,100 ... 2,000 & \centering о.е. & \centering 0,001 & \centering \arraybackslash 0,200 \\
\hline
\raggedright Iт1 & \centering 0,200 ... 1,500 & \centering о.е. & \centering 0,001 & \centering \arraybackslash 1,000 \\
\hline
\raggedright Кт1 & \centering 0,100 ... 1,000 & \centering о.е. & \centering 0,001 & \centering \arraybackslash 0,500 \\
\hline
\raggedright Iт2 & \centering 1,000 ... 3,000 & \centering о.е. & \centering 0,001 & \centering \arraybackslash 2,000 \\
\hline
\raggedright Кт2 & \centering 0,200 ... 1,000 & \centering о.е. & \centering 0,001 & \centering \arraybackslash 0,750 \\
\hline
\raggedright Кв (Iср диф) & \centering 0,00 ... 1,00 & \centering  & \centering 0,01 & \centering \arraybackslash 0,80 \\
\hline
\raggedright фВ & \centering Введено \\ Выведено & \centering --- & \centering --- & \centering \arraybackslash Введено \\
\hline
\raggedright ДТЗт & \centering Введено \\ Выведено & \centering --- & \centering --- & \centering \arraybackslash Введено \\
\hline
\multicolumn{5}{|c|}{ ДТО } \\ \hline 
\raggedright Iср ДТО & \centering 1,000 ... 20,000 & \centering о.е. & \centering 0,001 & \centering \arraybackslash 20,000 \\
\hline
\raggedright Кв (Iср) & \centering 0,00 ... 1,00 & \centering  & \centering 0,01 & \centering \arraybackslash 0,95 \\
\hline
\raggedright фВ & \centering Введено \\ Выведено & \centering --- & \centering --- & \centering \arraybackslash Введено \\
\hline
\raggedright ДТО & \centering Введено \\ Выведено & \centering --- & \centering --- & \centering \arraybackslash Введено \\
\hline
\multicolumn{5}{|c|}{ Д2г } \\ \hline 
\raggedright I2г/I1г ДТЗт & \centering 0,050 ... 1,000 & \centering о.е. & \centering 0,001 & \centering \arraybackslash 0,500 \\
\hline
\raggedright Кв (I2(5)г/I1г) & \centering 0,00 ... 1,00 & \centering  & \centering 0,01 & \centering \arraybackslash 0,95 \\
\hline
\raggedright Iср диф & \centering 0,000 ... 2,000 & \centering о.е. & \centering 0,001 & \centering \arraybackslash 0,200 \\
\hline
\raggedright Кв (Iср диф) & \centering 0,00 ... 1,00 & \centering  & \centering 0,01 & \centering \arraybackslash 0,95 \\
\hline
\raggedright 0.01 с & \centering 0 ... 10 & \centering мc & \centering 5 & \centering \arraybackslash 10 \\
\hline
\raggedright 0.02 с & \centering 0 ... 20 & \centering мc & \centering 5 & \centering \arraybackslash 20 \\
\hline
\raggedright фВ & \centering Введено \\ Выведено & \centering --- & \centering --- & \centering \arraybackslash Введено \\
\hline
\raggedright Перекр блок 2г & \centering Откл \\ 1 из 3 \\ 2 из 3 & \centering --- & \centering --- & \centering \arraybackslash Откл \\
\hline
\raggedright Тперекр блок 2г & \centering 100 ... 1000 & \centering мc & \centering 5 & \centering \arraybackslash 100 \\
\hline
\raggedright Блок 2г ДТЗт & \centering Введено \\ Выведено & \centering --- & \centering --- & \centering \arraybackslash Введено \\
\hline
\multicolumn{5}{|c|}{ Д5г } \\ \hline 
\raggedright I5г/I1г ДТЗт & \centering 0,050 ... 1,000 & \centering о.е. & \centering 0,001 & \centering \arraybackslash 0,500 \\
\hline
\raggedright Кв (I2(5)г/I1г) & \centering 0,00 ... 1,00 & \centering  & \centering 0,01 & \centering \arraybackslash 0,95 \\
\hline
\raggedright Iср диф & \centering 0,000 ... 2,000 & \centering о.е. & \centering 0,001 & \centering \arraybackslash 0,200 \\
\hline
\raggedright Кв (Iср диф) & \centering 0,00 ... 1,00 & \centering  & \centering 0,01 & \centering \arraybackslash 0,95 \\
\hline
\raggedright 0.01 с & \centering 0 ... 10 & \centering мc & \centering 5 & \centering \arraybackslash 10 \\
\hline
\raggedright 0.02 с & \centering 0 ... 20 & \centering мc & \centering 5 & \centering \arraybackslash 20 \\
\hline
\raggedright фВ & \centering Введено \\ Выведено & \centering --- & \centering --- & \centering \arraybackslash Введено \\
\hline
\raggedright Перекр блок 5г & \centering Откл \\ 1 из 3 \\ 2 из 3 & \centering --- & \centering --- & \centering \arraybackslash Откл \\
\hline
\raggedright Тперекр блок 5г & \centering 100 ... 1000 & \centering мc & \centering 5 & \centering \arraybackslash 100 \\
\hline
\raggedright Блок 5г ДТЗт & \centering Введено \\ Выведено & \centering --- & \centering --- & \centering \arraybackslash Введено \\
\hline
%===m
\stopTable


%%%%%%%%%%%%%%%%%%%%%%%%%%%%%%%%%%%%%%% МТЗ НН

\par
Уставки МТЗ НН представлены в таблице \ref{app_set:mtz}.

\startTable{Параметры для настройки МТЗ НН}{app_set:mtz}

%===m>МТЗ НН|general+-
\multicolumn{5}{|c|}{ МТЗ-1 } \\ \hline 
\raggedright Iср МТЗ-1 НН & \centering 0,04 ... 30,00 & \centering о.е. & \centering 0,01 & \centering \arraybackslash 30,00 \\
\hline
\raggedright Тср МТЗ-1 НН & \centering 0 ... 300000 & \centering мс & \centering 5 & \centering \arraybackslash 300000 \\
\hline
\raggedright МТЗ-1 НН & \centering Введено \\ Выведено & \centering --- & \centering --- & \centering \arraybackslash Введено \\
\hline
\raggedright Режим КПОН МТЗ-1 НН & \centering Введено \\ Выведено & \centering --- & \centering --- & \centering \arraybackslash Введено \\
\hline
\raggedright Режим БНТ МТЗ-1 НН & \centering Введено \\ Выведено & \centering --- & \centering --- & \centering \arraybackslash Введено \\
\hline
\multicolumn{5}{|c|}{ МТЗ-2 } \\ \hline 
\raggedright Iср МТЗ-2 НН & \centering 0,04 ... 30,00 & \centering о.е. & \centering 0,01 & \centering \arraybackslash 30,00 \\
\hline
\raggedright Тср МТЗ-2 НН & \centering 0 ... 300000 & \centering мс & \centering 5 & \centering \arraybackslash 300000 \\
\hline
\raggedright МТЗ-2 НН & \centering Введено \\ Выведено & \centering --- & \centering --- & \centering \arraybackslash Введено \\
\hline
\raggedright Режим КПОН МТЗ-2 НН & \centering Введено \\ Выведено & \centering --- & \centering --- & \centering \arraybackslash Введено \\
\hline
\raggedright Режим БНТ МТЗ-2 НН & \centering Введено \\ Выведено & \centering --- & \centering --- & \centering \arraybackslash Введено \\
\hline
\multicolumn{5}{|c|}{ ТО } \\ \hline 
\raggedright Iср ТО НН & \centering 0,04 ... 30,00 & \centering о.е. & \centering 0,01 & \centering \arraybackslash 30,00 \\
\hline
\raggedright Тср ТО НН & \centering 0 ... 300000 & \centering мс & \centering 5 & \centering \arraybackslash 300000 \\
\hline
\raggedright  ТО НН & \centering Введено \\ Выведено & \centering --- & \centering --- & \centering \arraybackslash Введено \\
\hline
\raggedright Режим БНТ ТО НН & \centering Введено \\ Выведено & \centering --- & \centering --- & \centering \arraybackslash Введено \\
\hline
\multicolumn{5}{|c|}{ БНТ } \\ \hline 
\raggedright I2г/I1г МТЗ НН & \centering 0,00 ... 1,00 & \centering о.е. & \centering 0,01 & \centering \arraybackslash 0,20 \\
\hline
\raggedright kв=0.95 & \centering 0,00 ... 1,00 & \centering  & \centering 0,01 & \centering \arraybackslash 0,95 \\
\hline
\raggedright 0.04Iном (Iразр БНТ МТЗ НН) & \centering 0,04 ... 30,00 & \centering о.е & \centering 0,01 & \centering \arraybackslash 0,04 \\
\hline
\raggedright kв=1 & \centering 0,00 ... 1,00 & \centering  & \centering 0,01 & \centering \arraybackslash 1,00 \\
\hline
\raggedright БНТ МТЗ НН & \centering Введено \\ Выведено & \centering --- & \centering --- & \centering \arraybackslash Введено \\
\hline
\raggedright Iблок БНТ МТЗ НН & \centering 0,04 ... 30,00 & \centering о.е. & \centering 0,01 & \centering \arraybackslash 30,00 \\
\hline
%===m
\end{longtable}
\normalsize


%%%%%%%%%%%%%%%%%%%%%%%%%%%%%%%%%%%%%%% КПОН

\par
Уставки КПОН представлены в таблице \ref{app_set:kpon}.

\startTable{Параметры для настройки КПОН}{app_set:kpon}

%===m>КПОН|general+-
\multicolumn{5}{|c|}{ КПОН НН1 } \\ \hline 
\raggedright U2 КПОН НН1 & \centering 0,0050 ... 1,5000 & \centering о.е & \centering 0,0001 & \centering \arraybackslash 1,5000 \\
\hline
\raggedright Uл КПОН НН1 & \centering 0,0050 ... 1,5000 & \centering о.е & \centering 0,0001 & \centering \arraybackslash 0,0050 \\
\hline
\raggedright Тип КПОН НН1 & \centering Umin, U2 \\ Umin \\ Внеш & \centering --- & \centering --- & \centering \arraybackslash Umin, U2 \\
\hline
\raggedright  КПОН НН1 & \centering Введено \\ Выведено & \centering --- & \centering --- & \centering \arraybackslash Введено \\
\hline
\raggedright Неиспр ТН КПОН НН1 & \centering Выведено \\ Пуск КПОН \\ Блок КПОН & \centering --- & \centering --- & \centering \arraybackslash Выведено \\
\hline
\multicolumn{5}{|c|}{ КПОН НН2 } \\ \hline 
\raggedright U2 КПОН НН2 & \centering 0,0050 ... 1,5000 & \centering о.е & \centering 0,0001 & \centering \arraybackslash 1,5000 \\
\hline
\raggedright Uл КПОН НН2 & \centering 0,0050 ... 1,5000 & \centering о.е & \centering 0,0001 & \centering \arraybackslash 0,0050 \\
\hline
\raggedright Тип КПОН НН2 & \centering Umin, U2 \\ Umin \\ Внеш & \centering --- & \centering --- & \centering \arraybackslash Umin, U2 \\
\hline
\raggedright  КПОН НН2 & \centering Введено \\ Выведено & \centering --- & \centering --- & \centering \arraybackslash Введено \\
\hline
\raggedright Неиспр ТН КПОН НН2 & \centering Выведено \\ Пуск КПОН \\ Блок КПОН & \centering --- & \centering --- & \centering \arraybackslash Выведено \\
\hline
\multicolumn{5}{|c|}{ КПОН НН } \\ \hline 
\raggedright U2 КПОН НН & \centering 0,0050 ... 1,5000 & \centering о.е & \centering 0,0001 & \centering \arraybackslash 1,5000 \\
\hline
\raggedright Uл КПОН НН & \centering 0,0050 ... 1,5000 & \centering о.е & \centering 0,0001 & \centering \arraybackslash 0,0050 \\
\hline
\raggedright Тип КПОН НН & \centering Umin, U2 \\ Umin \\ Внеш & \centering --- & \centering --- & \centering \arraybackslash Umin, U2 \\
\hline
\raggedright КПОН НН & \centering Введено \\ Выведено & \centering --- & \centering --- & \centering \arraybackslash Введено \\
\hline
\raggedright Неиспр ТН КПОН НН & \centering Выведено \\ Пуск КПОН \\ Блок КПОН & \centering --- & \centering --- & \centering \arraybackslash Выведено \\
\hline
%===m
\end{longtable}
\normalsize



%%%%%%%%%%%%%%%%%%%%%%%%%%%%%%%%%%%%%%% ГЗ АТ откл

\par
Уставки отключающей ступени ГЗ АТ представлены в таблице \ref{app_set:gzotkl}.

\startTable{Параметры для настройки отключающей ступени ГЗ АТ}{app_set:gzotkl}

%===m>ГЗ АТ откл|general+-
\multicolumn{5}{|c|}{ ГЗ АТ откл } \\ \hline 
\raggedright ГЗ АТ откл 1 цепь & \centering Введено \\ Выведено & \centering --- & \centering --- & \centering \arraybackslash Введено \\
\hline
\raggedright ГЗ АТ откл 2 цепь & \centering Введено \\ Выведено & \centering --- & \centering --- & \centering \arraybackslash Введено \\
\hline
\raggedright Блокировка ГЗ АТ откл & \centering Введено \\ Выведено & \centering --- & \centering --- & \centering \arraybackslash Введено \\
\hline
\raggedright Тбл ГЗ АТ откл & \centering 0 ... 300000 & \centering мс & \centering 5 & \centering \arraybackslash 300000 \\
\hline
\multicolumn{5}{|c|}{ ГЗ откл 1 цепь } \\ \hline 
\raggedright Блокировка & \centering Введено \\ Выведено & \centering --- & \centering --- & \centering \arraybackslash Введено \\
\hline
\raggedright ГЗ (ТЗ) откл & \centering Введено \\ Выведено & \centering --- & \centering --- & \centering \arraybackslash Введено \\
\hline
\raggedright Тбл & \centering 0 ... 300000 & \centering мс & \centering 5 & \centering \arraybackslash 300000 \\
\hline
\multicolumn{5}{|c|}{ ГЗ откл 2 цепь } \\ \hline 
\raggedright Блокировка & \centering Введено \\ Выведено & \centering --- & \centering --- & \centering \arraybackslash Введено \\
\hline
\raggedright ГЗ (ТЗ) откл & \centering Введено \\ Выведено & \centering --- & \centering --- & \centering \arraybackslash Введено \\
\hline
\raggedright Тбл & \centering 0 ... 300000 & \centering мс & \centering 5 & \centering \arraybackslash 300000 \\
\hline
%===m
\end{longtable}
\normalsize


%%%%%%%%%%%%%%%%%%%%%%%%%%%%%%%%%%%%%%% ГЗ АТ сигн

\par
Уставки сигнальной ступени ГЗ АТ представлены в таблице \ref{app_set:gzsign}.

\startTable{Параметры для настройки сигнальной ступени ГЗ АТ}{app_set:gzsign}

%===m>ГЗ АТ сигн|general+-
\multicolumn{5}{|c|}{ ГЗ АТ сигн } \\ \hline 
\raggedright ГЗ АТ сигн 1 цепь & \centering Введено \\ Выведено & \centering --- & \centering --- & \centering \arraybackslash Введено \\
\hline
\raggedright ГЗ АТ сигн 2 цепь & \centering Введено \\ Выведено & \centering --- & \centering --- & \centering \arraybackslash Введено \\
\hline
\raggedright Откл от ГЗ АТ сигн & \centering Введено \\ Выведено & \centering --- & \centering --- & \centering \arraybackslash Введено \\
\hline
\multicolumn{5}{|c|}{ ГЗ сигн 2 цепь } \\ \hline 
\raggedright ГЗ (ТЗ) сигн & \centering Введено \\ Выведено & \centering --- & \centering --- & \centering \arraybackslash Введено \\
\hline
\multicolumn{5}{|c|}{ ГЗ сигн 1 цепь } \\ \hline 
\raggedright ГЗ (ТЗ) сигн & \centering Введено \\ Выведено & \centering --- & \centering --- & \centering \arraybackslash Введено \\
\hline
%===m
\end{longtable}
\normalsize


%%%%%%%%%%%%%%%%%%%%%%%%%%%%%%%%%%%%%%% ГЗ ЛРТ откл

\par
Уставки отключающей ступени ГЗ ЛРТ представлены в таблице \ref{app_set:gzlrtotkl}.

\startTable{Параметры для настройки отключающей ступени ГЗ ЛРТ}{app_set:gzlrtotkl}

%===m>ГЗ ЛРТ откл|general+-
\multicolumn{5}{|c|}{ ГЗ ЛРТ откл } \\ \hline 
\raggedright ГЗ ЛРТ откл 1 цепь & \centering Введено \\ Выведено & \centering --- & \centering --- & \centering \arraybackslash Введено \\
\hline
\raggedright ГЗ ЛРТ откл 2 цепь & \centering Введено \\ Выведено & \centering --- & \centering --- & \centering \arraybackslash Введено \\
\hline
\raggedright Блокировка ГЗ ЛРТ & \centering Введено \\ Выведено & \centering --- & \centering --- & \centering \arraybackslash Введено \\
\hline
\raggedright Тбл ГЗ ЛРТ & \centering 0 ... 300000 & \centering мс & \centering 5 & \centering \arraybackslash 300000 \\
\hline
\multicolumn{5}{|c|}{ ГЗ откл 1 цепь } \\ \hline 
\raggedright Блокировка & \centering Введено \\ Выведено & \centering --- & \centering --- & \centering \arraybackslash Введено \\
\hline
\raggedright ГЗ (ТЗ) откл & \centering Введено \\ Выведено & \centering --- & \centering --- & \centering \arraybackslash Введено \\
\hline
\raggedright Тбл & \centering 0 ... 300000 & \centering мс & \centering 5 & \centering \arraybackslash 300000 \\
\hline
\multicolumn{5}{|c|}{ ГЗ откл 2 цепь } \\ \hline 
\raggedright Блокировка & \centering Введено \\ Выведено & \centering --- & \centering --- & \centering \arraybackslash Введено \\
\hline
\raggedright ГЗ (ТЗ) откл & \centering Введено \\ Выведено & \centering --- & \centering --- & \centering \arraybackslash Введено \\
\hline
\raggedright Тбл & \centering 0 ... 300000 & \centering мс & \centering 5 & \centering \arraybackslash 300000 \\
\hline
%===m
\end{longtable}
\normalsize


%%%%%%%%%%%%%%%%%%%%%%%%%%%%%%%%%%%%%%% ГЗ ЛРТ сигн

\par
Уставки сигнальной ступени ГЗ ЛРТ представлены в таблице \ref{app_set:gzlrtsign}.

\startTable{Параметры для настройки сигнальной ступени ГЗ ЛРТ}{app_set:gzlrtsign}

%===m>ГЗ ЛРТ сигн|general+-
\multicolumn{5}{|c|}{ ГЗ ЛРТ сигн } \\ \hline 
\raggedright ГЗ ЛРТ сигн 1 цепь & \centering Введено \\ Выведено & \centering --- & \centering --- & \centering \arraybackslash Введено \\
\hline
\raggedright ГЗ ЛРТ сигн 2 цепь & \centering Введено \\ Выведено & \centering --- & \centering --- & \centering \arraybackslash Введено \\
\hline
\raggedright Откл от ГЗ ЛРТ сигн & \centering Введено \\ Выведено & \centering --- & \centering --- & \centering \arraybackslash Введено \\
\hline
\multicolumn{5}{|c|}{ ГЗ сигн 2 цепь } \\ \hline 
\raggedright ГЗ (ТЗ) сигн & \centering Введено \\ Выведено & \centering --- & \centering --- & \centering \arraybackslash Введено \\
\hline
\multicolumn{5}{|c|}{ ГЗ сигн 1 цепь } \\ \hline 
\raggedright ГЗ (ТЗ) сигн & \centering Введено \\ Выведено & \centering --- & \centering --- & \centering \arraybackslash Введено \\
\hline
%===m
\end{longtable}
\normalsize


%%%%%%%%%%%%%%%%%%%%%%%%%%%%%%%%%%%%%%% ГЗ РПН

\par
Уставки ГЗ РПН представлены в таблице \ref{app_set:gzrpn}.

\startTable{Параметры для настройки ГЗ РПН}{app_set:gzrpn}

%===m>ГЗ РПН|general+-
\multicolumn{5}{|c|}{ ГЗ РПН } \\ \hline 
\raggedright ГЗ РПН 1 цепь & \centering Введено \\ Выведено & \centering --- & \centering --- & \centering \arraybackslash Введено \\
\hline
\raggedright ГЗ РПН 2 цепь & \centering Введено \\ Выведено & \centering --- & \centering --- & \centering \arraybackslash Введено \\
\hline
\raggedright Блок ГЗ РПН & \centering Введено \\ Выведено & \centering --- & \centering --- & \centering \arraybackslash Введено \\
\hline
\raggedright Тбл ГЗ РПН & \centering 0 ... 300000 & \centering мс & \centering 5 & \centering \arraybackslash 300000 \\
\hline
\multicolumn{5}{|c|}{ ГЗ откл.фА 1 цепь } \\ \hline 
\raggedright Блокировка & \centering Введено \\ Выведено & \centering --- & \centering --- & \centering \arraybackslash Введено \\
\hline
\raggedright ГЗ (ТЗ) откл & \centering Введено \\ Выведено & \centering --- & \centering --- & \centering \arraybackslash Введено \\
\hline
\raggedright Тбл & \centering 0 ... 300000 & \centering мс & \centering 5 & \centering \arraybackslash 300000 \\
\hline
\multicolumn{5}{|c|}{ ГЗ откл.фА 2 цепь } \\ \hline 
\raggedright Блокировка & \centering Введено \\ Выведено & \centering --- & \centering --- & \centering \arraybackslash Введено \\
\hline
\raggedright ГЗ (ТЗ) откл & \centering Введено \\ Выведено & \centering --- & \centering --- & \centering \arraybackslash Введено \\
\hline
\raggedright Тбл & \centering 0 ... 300000 & \centering мс & \centering 5 & \centering \arraybackslash 300000 \\
\hline
\multicolumn{5}{|c|}{ ГЗ откл.фВ 1 цепь } \\ \hline 
\raggedright Блокировка & \centering Введено \\ Выведено & \centering --- & \centering --- & \centering \arraybackslash Введено \\
\hline
\raggedright ГЗ (ТЗ) откл & \centering Введено \\ Выведено & \centering --- & \centering --- & \centering \arraybackslash Введено \\
\hline
\raggedright Тбл & \centering 0 ... 300000 & \centering мс & \centering 5 & \centering \arraybackslash 300000 \\
\hline
\multicolumn{5}{|c|}{ ГЗ откл.фВ 2 цепь } \\ \hline 
\raggedright Блокировка & \centering Введено \\ Выведено & \centering --- & \centering --- & \centering \arraybackslash Введено \\
\hline
\raggedright ГЗ (ТЗ) откл & \centering Введено \\ Выведено & \centering --- & \centering --- & \centering \arraybackslash Введено \\
\hline
\raggedright Тбл & \centering 0 ... 300000 & \centering мс & \centering 5 & \centering \arraybackslash 300000 \\
\hline
\multicolumn{5}{|c|}{ ГЗ откл.фС 1 цепь } \\ \hline 
\raggedright Блокировка & \centering Введено \\ Выведено & \centering --- & \centering --- & \centering \arraybackslash Введено \\
\hline
\raggedright ГЗ (ТЗ) откл & \centering Введено \\ Выведено & \centering --- & \centering --- & \centering \arraybackslash Введено \\
\hline
\raggedright Тбл & \centering 0 ... 300000 & \centering мс & \centering 5 & \centering \arraybackslash 300000 \\
\hline
\multicolumn{5}{|c|}{ ГЗ откл.фС 2 цепь } \\ \hline 
\raggedright Блокировка & \centering Введено \\ Выведено & \centering --- & \centering --- & \centering \arraybackslash Введено \\
\hline
\raggedright ГЗ (ТЗ) откл & \centering Введено \\ Выведено & \centering --- & \centering --- & \centering \arraybackslash Введено \\
\hline
\raggedright Тбл & \centering 0 ... 300000 & \centering мс & \centering 5 & \centering \arraybackslash 300000 \\
\hline
%===m
\end{longtable}
\normalsize


%%%%%%%%%%%%%%%%%%%%%%%%%%%%%%%%%%%%%%% ТЗ АТ

\par
Уставки ТЗ АТ представлены в таблице \ref{app_set:tzat}.

\startTable{Параметры для настройки ТЗ АТ}{app_set:tzat}

%===m>ТЗ АТ|general+-
\multicolumn{5}{|c|}{ ТЗ АТ } \\ \hline 
\raggedright ДТм АТ сигн 1 цепь & \centering Введено \\ Выведено & \centering --- & \centering --- & \centering \arraybackslash Введено \\
\hline
\raggedright ДТм АТ сигн 2 цепь & \centering Введено \\ Выведено & \centering --- & \centering --- & \centering \arraybackslash Введено \\
\hline
\raggedright ДТм АТ авар 1 цепь & \centering Введено \\ Выведено & \centering --- & \centering --- & \centering \arraybackslash Введено \\
\hline
\raggedright ДТм АТ авар 2 цепь & \centering Введено \\ Выведено & \centering --- & \centering --- & \centering \arraybackslash Введено \\
\hline
\raggedright ДТо АТ сигн 1 цепь & \centering Введено \\ Выведено & \centering --- & \centering --- & \centering \arraybackslash Введено \\
\hline
\raggedright ДТо АТ сигн 2 цепь & \centering Введено \\ Выведено & \centering --- & \centering --- & \centering \arraybackslash Введено \\
\hline
\raggedright ДТо АТ авар 1 цепь & \centering Введено \\ Выведено & \centering --- & \centering --- & \centering \arraybackslash Введено \\
\hline
\raggedright ДТо АТ авар 2 цепь & \centering Введено \\ Выведено & \centering --- & \centering --- & \centering \arraybackslash Введено \\
\hline
\raggedright Блокировка ТЗ АТ & \centering Введено \\ Выведено & \centering --- & \centering --- & \centering \arraybackslash Введено \\
\hline
\raggedright Тбл ТЗ АТ & \centering 0 ... 300000 & \centering мс & \centering 5 & \centering \arraybackslash 300000 \\
\hline
\raggedright Откл от ДТм АТ авар & \centering Введено \\ Выведено & \centering --- & \centering --- & \centering \arraybackslash Введено \\
\hline
\raggedright Откл от ДТо АТ авар & \centering Введено \\ Выведено & \centering --- & \centering --- & \centering \arraybackslash Введено \\
\hline
\raggedright Откл от отсеч клап АТ & \centering Введено \\ Выведено & \centering --- & \centering --- & \centering \arraybackslash Введено \\
\hline
\raggedright Откл от предохр клап АТ & \centering Введено \\ Выведено & \centering --- & \centering --- & \centering \arraybackslash Введено \\
\hline
\multicolumn{5}{|c|}{ ДТм сигн 1 цепь } \\ \hline 
\raggedright ГЗ (ТЗ) сигн & \centering Введено \\ Выведено & \centering --- & \centering --- & \centering \arraybackslash Введено \\
\hline
\multicolumn{5}{|c|}{ ДТм сигн 2 цепь } \\ \hline 
\raggedright ГЗ (ТЗ) сигн & \centering Введено \\ Выведено & \centering --- & \centering --- & \centering \arraybackslash Введено \\
\hline
\multicolumn{5}{|c|}{ ДТм авар 1 цепь } \\ \hline 
\raggedright Блокировка & \centering Введено \\ Выведено & \centering --- & \centering --- & \centering \arraybackslash Введено \\
\hline
\raggedright ГЗ (ТЗ) откл & \centering Введено \\ Выведено & \centering --- & \centering --- & \centering \arraybackslash Введено \\
\hline
\raggedright Тбл & \centering 0 ... 300000 & \centering мс & \centering 5 & \centering \arraybackslash 300000 \\
\hline
\multicolumn{5}{|c|}{ ДТм авар 2 цепь } \\ \hline 
\raggedright Блокировка & \centering Введено \\ Выведено & \centering --- & \centering --- & \centering \arraybackslash Введено \\
\hline
\raggedright ГЗ (ТЗ) откл & \centering Введено \\ Выведено & \centering --- & \centering --- & \centering \arraybackslash Введено \\
\hline
\raggedright Тбл & \centering 0 ... 300000 & \centering мс & \centering 5 & \centering \arraybackslash 300000 \\
\hline
\multicolumn{5}{|c|}{ ДТо сигн 1 цепь } \\ \hline 
\raggedright ГЗ (ТЗ) сигн & \centering Введено \\ Выведено & \centering --- & \centering --- & \centering \arraybackslash Введено \\
\hline
\multicolumn{5}{|c|}{ ДТо сигн 2 цепь } \\ \hline 
\raggedright ГЗ (ТЗ) сигн & \centering Введено \\ Выведено & \centering --- & \centering --- & \centering \arraybackslash Введено \\
\hline
\multicolumn{5}{|c|}{ ДТо авар 1 цепь } \\ \hline 
\raggedright Блокировка & \centering Введено \\ Выведено & \centering --- & \centering --- & \centering \arraybackslash Введено \\
\hline
\raggedright ГЗ (ТЗ) откл & \centering Введено \\ Выведено & \centering --- & \centering --- & \centering \arraybackslash Введено \\
\hline
\raggedright Тбл & \centering 0 ... 300000 & \centering мс & \centering 5 & \centering \arraybackslash 300000 \\
\hline
\multicolumn{5}{|c|}{ ДТо авар 2 цепь } \\ \hline 
\raggedright Блокировка & \centering Введено \\ Выведено & \centering --- & \centering --- & \centering \arraybackslash Введено \\
\hline
\raggedright ГЗ (ТЗ) откл & \centering Введено \\ Выведено & \centering --- & \centering --- & \centering \arraybackslash Введено \\
\hline
\raggedright Тбл & \centering 0 ... 300000 & \centering мс & \centering 5 & \centering \arraybackslash 300000 \\
\hline
%===m
\end{longtable}
\normalsize


%%%%%%%%%%%%%%%%%%%%%%%%%%%%%%%%%%%%%%% ТЗ ЛРТ

\par
Уставки ТЗ ЛРТ представлены в таблице \ref{app_set:tzlrt}.

\startTable{Параметры для настройки ТЗ ЛРТ}{app_set:tzlrt}

%===m>ТЗ ЛРТ|general+-
\multicolumn{5}{|c|}{ ТЗ ЛРТ } \\ \hline 
\raggedright ДТм ЛРТ сигн 1 цепь & \centering Введено \\ Выведено & \centering --- & \centering --- & \centering \arraybackslash Введено \\
\hline
\raggedright ДТм ЛРТ сигн 2 цепь & \centering Введено \\ Выведено & \centering --- & \centering --- & \centering \arraybackslash Введено \\
\hline
\raggedright ДТм ЛРТ авар 1 цепь & \centering Введено \\ Выведено & \centering --- & \centering --- & \centering \arraybackslash Введено \\
\hline
\raggedright ДТм ЛРТ авар 2 цепь & \centering Введено \\ Выведено & \centering --- & \centering --- & \centering \arraybackslash Введено \\
\hline
\raggedright ДТо ЛРТ сигн 1 цепь & \centering Введено \\ Выведено & \centering --- & \centering --- & \centering \arraybackslash Введено \\
\hline
\raggedright ДТо ЛРТсигн 2 цепь & \centering Введено \\ Выведено & \centering --- & \centering --- & \centering \arraybackslash Введено \\
\hline
\raggedright ДТо ЛРТ авар 1 цепь & \centering Введено \\ Выведено & \centering --- & \centering --- & \centering \arraybackslash Введено \\
\hline
\raggedright ДТо ЛРТ авар 2 цепь & \centering Введено \\ Выведено & \centering --- & \centering --- & \centering \arraybackslash Введено \\
\hline
\raggedright РД ЛРТ 2 цепь & \centering Введено \\ Выведено & \centering --- & \centering --- & \centering \arraybackslash Введено \\
\hline
\raggedright РД ЛРТ 2 цепь & \centering Введено \\ Выведено & \centering --- & \centering --- & \centering \arraybackslash Введено \\
\hline
\raggedright Блокировка ТЗ ЛРТ & \centering Введено \\ Выведено & \centering --- & \centering --- & \centering \arraybackslash Введено \\
\hline
\raggedright Тбл ТЗ ЛРТ & \centering 0 ... 300000 & \centering мс & \centering 5 & \centering \arraybackslash 300000 \\
\hline
\raggedright Откл от ДТм ЛРТ авар & \centering Введено \\ Выведено & \centering --- & \centering --- & \centering \arraybackslash Введено \\
\hline
\raggedright Откл от ДТо  ЛРТ авар & \centering Введено \\ Выведено & \centering --- & \centering --- & \centering \arraybackslash Введено \\
\hline
\raggedright Откл от РД ЛРТ & \centering Введено \\ Выведено & \centering --- & \centering --- & \centering \arraybackslash Введено \\
\hline
\raggedright Откл от предохр клап ЛРТ & \centering Введено \\ Выведено & \centering --- & \centering --- & \centering \arraybackslash Выведено \\
\hline
\multicolumn{5}{|c|}{ ДТм сигн 1 цепь } \\ \hline 
\raggedright ГЗ (ТЗ) сигн & \centering Введено \\ Выведено & \centering --- & \centering --- & \centering \arraybackslash Введено \\
\hline
\multicolumn{5}{|c|}{ ДТм сигн 2 цепь } \\ \hline 
\raggedright ГЗ (ТЗ) сигн & \centering Введено \\ Выведено & \centering --- & \centering --- & \centering \arraybackslash Введено \\
\hline
\multicolumn{5}{|c|}{ ДТо сигн 1 цепь } \\ \hline 
\raggedright ГЗ (ТЗ) сигн & \centering Введено \\ Выведено & \centering --- & \centering --- & \centering \arraybackslash Введено \\
\hline
\multicolumn{5}{|c|}{ ДТо сигн 2 цепь } \\ \hline 
\raggedright ГЗ (ТЗ) сигн & \centering Введено \\ Выведено & \centering --- & \centering --- & \centering \arraybackslash Введено \\
\hline
\multicolumn{5}{|c|}{ ДТм авар 1 цепь } \\ \hline 
\raggedright Блокировка & \centering Введено \\ Выведено & \centering --- & \centering --- & \centering \arraybackslash Введено \\
\hline
\raggedright ГЗ (ТЗ) откл & \centering Введено \\ Выведено & \centering --- & \centering --- & \centering \arraybackslash Введено \\
\hline
\raggedright Тбл & \centering 0 ... 300000 & \centering мс & \centering 5 & \centering \arraybackslash 300000 \\
\hline
\multicolumn{5}{|c|}{ ДТм авар 2 цепь } \\ \hline 
\raggedright Блокировка & \centering Введено \\ Выведено & \centering --- & \centering --- & \centering \arraybackslash Введено \\
\hline
\raggedright ГЗ (ТЗ) откл & \centering Введено \\ Выведено & \centering --- & \centering --- & \centering \arraybackslash Введено \\
\hline
\raggedright Тбл & \centering 0 ... 300000 & \centering мс & \centering 5 & \centering \arraybackslash 300000 \\
\hline
\multicolumn{5}{|c|}{ ДТо авар 1 цепь } \\ \hline 
\raggedright Блокировка & \centering Введено \\ Выведено & \centering --- & \centering --- & \centering \arraybackslash Введено \\
\hline
\raggedright ГЗ (ТЗ) откл & \centering Введено \\ Выведено & \centering --- & \centering --- & \centering \arraybackslash Введено \\
\hline
\raggedright Тбл & \centering 0 ... 300000 & \centering мс & \centering 5 & \centering \arraybackslash 300000 \\
\hline
\multicolumn{5}{|c|}{ ДТо авар 2 цепь } \\ \hline 
\raggedright Блокировка & \centering Введено \\ Выведено & \centering --- & \centering --- & \centering \arraybackslash Введено \\
\hline
\raggedright ГЗ (ТЗ) откл & \centering Введено \\ Выведено & \centering --- & \centering --- & \centering \arraybackslash Введено \\
\hline
\raggedright Тбл & \centering 0 ... 300000 & \centering мс & \centering 5 & \centering \arraybackslash 300000 \\
\hline
\multicolumn{5}{|c|}{ РД ЛРТ 1 цепь } \\ \hline 
\raggedright Блокировка & \centering Введено \\ Выведено & \centering --- & \centering --- & \centering \arraybackslash Введено \\
\hline
\raggedright ГЗ (ТЗ) откл & \centering Введено \\ Выведено & \centering --- & \centering --- & \centering \arraybackslash Введено \\
\hline
\raggedright Тбл & \centering 0 ... 300000 & \centering мс & \centering 5 & \centering \arraybackslash 300000 \\
\hline
\multicolumn{5}{|c|}{ РД ЛРТ 2 цепь } \\ \hline 
\raggedright Блокировка & \centering Введено \\ Выведено & \centering --- & \centering --- & \centering \arraybackslash Введено \\
\hline
\raggedright ГЗ (ТЗ) откл & \centering Введено \\ Выведено & \centering --- & \centering --- & \centering \arraybackslash Введено \\
\hline
\raggedright Тбл & \centering 0 ... 300000 & \centering мс & \centering 5 & \centering \arraybackslash 300000 \\
\hline
%===m
\end{longtable}
\normalsize


%%%%%%%%%%%%%%%%%%%%%%%%%%%%%%%%%%%%%%% ТК ЗДЗ
\par
Уставки ТК ЗДЗ представлены в таблице \ref{app_set:tkzdz}.

\startTable{Параметры для настройки ТК ЗДЗ}{app_set:tkzdz}

%===m>ТК ЗДЗ|general+-
\multicolumn{5}{|c|}{ ТК ЗДЗ НН } \\ \hline 
\raggedright Iср ТК ЗДЗ НН & \centering 0,04 ... 30,00 & \centering о.е. & \centering 0,01 & \centering \arraybackslash 30,00 \\
\hline
\raggedright ТК ЗДЗ НН & \centering Введено \\ Выведено & \centering --- & \centering --- & \centering \arraybackslash Введено \\
\hline
\multicolumn{5}{|c|}{ ТК ЗДЗ СН } \\ \hline 
\raggedright Iср ТК ЗДЗ СН & \centering 0,04 ... 30,00 & \centering о.е. & \centering 0,01 & \centering \arraybackslash 30,00 \\
\hline
\raggedright ТК ЗДЗ СН & \centering Введено \\ Выведено & \centering --- & \centering --- & \centering \arraybackslash Введено \\
\hline
\multicolumn{5}{|c|}{ ТК ЗДЗ ВН } \\ \hline 
\raggedright Iср ТК ЗДЗ ВН & \centering 0,04 ... 30,00 & \centering о.е. & \centering 0,01 & \centering \arraybackslash 30,00 \\
\hline
\raggedright ТК ЗДЗ ВН & \centering Введено \\ Выведено & \centering --- & \centering --- & \centering \arraybackslash Введено \\
\hline
%===m
\end{longtable}
\normalsize


%%%%%%%%%%%%%%%%%%%%%%%%%%%%%%%%%%%%%%% АПТ

\par
Уставки логики пуска АПТ представлены в таблице \ref{app_set:apt}.

\startTable{Параметры для настройки логики пуска АПТ }{app_set:apt}

%===m>ЛППТ АТ|general+-
\multicolumn{5}{|c|}{ ЛППТ } \\ \hline 
\raggedright Тимп ЛППТ АТ  & \centering 0 ... 300000 & \centering мс & \centering 5 & \centering \arraybackslash 1000 \\
\hline
\raggedright ЛППТ АТ & \centering Введено \\ Выведено & \centering --- & \centering --- & \centering \arraybackslash Введено \\
\hline
%===m
\end{longtable}
\normalsize


%%%%%%%%%%%%%%%%%%%%%%%%%%%%%%%%%%%%%%% КОН

\par
Уставки КОН представлены в таблице \ref{app_set:kon}.

\startTable{Параметры для настройки КОН}{app_set:kon}

%===m>КОН|general+-
\multicolumn{5}{|c|}{ КОН } \\ \hline 
\raggedright КОН & \centering Введено \\ Выведено & \centering --- & \centering --- & \centering \arraybackslash Введено \\
\hline
\multicolumn{5}{|c|}{ РН НН } \\ \hline 
\raggedright Uср НН & \centering 0,0050 ... 1,5000 & \centering о.е. & \centering 0,0001 & \centering \arraybackslash 0,1000 \\
\hline
\raggedright РН НН & \centering Введено \\ Выведено & \centering --- & \centering --- & \centering \arraybackslash Введено \\
\hline
\multicolumn{5}{|c|}{ РТ ВН } \\ \hline 
\raggedright Iср РТ ВН & \centering 0,04 ... 40,00 & \centering о.е. & \centering 0,01 & \centering \arraybackslash 0,10 \\
\hline
\raggedright РТ ВН & \centering Введено \\ Выведено & \centering --- & \centering --- & \centering \arraybackslash Введено \\
\hline
\multicolumn{5}{|c|}{ РТ СН } \\ \hline 
\raggedright Iср РТ СН & \centering 0,04 ... 40,00 & \centering о.е. & \centering 0,01 & \centering \arraybackslash 0,10 \\
\hline
\raggedright РТ СН & \centering Введено \\ Выведено & \centering --- & \centering --- & \centering \arraybackslash Введено \\
\hline
%===m
\end{longtable}
\normalsize


%%%%%%%%%%%%%%%%%%%%%%%%%%%%%%%%%%%%%%% КОН

\par
Уставки РТПО представлены в таблице \ref{app_set:rtpo}.

\startTable{Параметры для настройки РТПО}{app_set:rtpo}


%===m>РТПО АТ|general+-
\multicolumn{5}{|c|}{ РТПО ВН } \\ \hline 
\raggedright Iср РТ1 ВН РТПО & \centering 0,04 ... 30,00 & \centering о.е. & \centering 0,01 & \centering \arraybackslash 30,00 \\
\hline
\raggedright Iср РТ2 ВН РТПО & \centering 0,04 ... 30,00 & \centering о.е. & \centering 0,01 & \centering \arraybackslash 30,00 \\
\hline
\raggedright РТ ВН РТПО & \centering Введено \\ Выведено & \centering --- & \centering --- & \centering \arraybackslash Введено \\
\hline
\multicolumn{5}{|c|}{ РТПО ОО } \\ \hline 
\raggedright Iср РТ1 ОО РТПО & \centering 0,04 ... 30,00 & \centering о.е. & \centering 0,01 & \centering \arraybackslash 30,00 \\
\hline
\raggedright Iср РТ2 ОО РТПО & \centering 0,04 ... 30,00 & \centering о.е. & \centering 0,01 & \centering \arraybackslash 30,00 \\
\hline
\raggedright РТ ОО РТПО & \centering Введено \\ Выведено & \centering --- & \centering --- & \centering \arraybackslash Введено \\
\hline
\multicolumn{5}{|c|}{ РТПО НН } \\ \hline 
\raggedright Iср РТ1 НН РТПО & \centering 0,04 ... 30,00 & \centering о.е. & \centering 0,01 & \centering \arraybackslash 30,00 \\
\hline
\raggedright Iср РТ2 НН РТПО & \centering 0,04 ... 30,00 & \centering о.е. & \centering 0,01 & \centering \arraybackslash 30,00 \\
\hline
\raggedright РТ НН РТПО & \centering Введено \\ Выведено & \centering --- & \centering --- & \centering \arraybackslash Введено \\
\hline
%===m
\end{longtable}
\normalsize



%%%%%%%%%%%%%%%%%%%%%%%%%%%%%%%%%%%%%%% РТ бл. РПН АТ

\par
Уставки РТ бл. РПН АТ представлены в таблице \ref{app_set:rtrpn}.

\footnotesize
\begin{longtable}{|>{\centering\arraybackslash}m{5.3cm}|>{\centering\arraybackslash}m{3.3cm}|>{\centering\arraybackslash}m{4.2cm}|>{\centering\arraybackslash}m{1.8cm}|>{\centering\arraybackslash}m{1cm}|}
% Первая страница: полный заголовок
\caption{Параметры для настройки РТ бл. РПН АТ\hfill\vspace{-0.5\baselineskip}}\label{app_set:rtrpn}\\
\hline
\rowcolor{gray!20}
Параметр & Условное обозначение на схеме & Значение/ Диапазон & Единица измерения & Шаг \\
\hline
\endfirsthead

% Промежуточные страницы: «Продолжение таблицы»
\caption*{\hspace{3pt}\emph{Продолжение таблицы \ref{app_set:rtrpn}\hfill\vspace{-0.5\baselineskip}}} \\ % сделано по ГОСТ 2.105 п.6.8.7

\hline
\rowcolor{gray!20}
Параметр & Условное обозначение на схеме & Значение/ Диапазон & Единица измерения & Шаг \\
\hline
\endhead

%===m>РТ бл. РПН АТ|general+-
\multicolumn{5}{|c|}{ РТ РПН } \\ \hline 
\raggedright Iбл РПН & \centering 0,04 ... 30,00 & \centering о.е & \centering 0,01 & \centering \arraybackslash 30,00 \\
\hline
\raggedright РТ бл РПН & \centering Введено \\ Выведено & \centering --- & \centering --- & \centering \arraybackslash Выведено \\
\hline
%===m
\end{longtable}
\normalsize


%%%%%%%%%%%%%%%%%%%%%%%%%%%%%%%%%%%%%%% ЗП
\par
Уставки ЗП представлены в таблице \ref{app_set:zp}.

\footnotesize
\begin{longtable}{|>{\centering\arraybackslash}m{5.3cm}|>{\centering\arraybackslash}m{3.3cm}|>{\centering\arraybackslash}m{4.2cm}|>{\centering\arraybackslash}m{1.8cm}|>{\centering\arraybackslash}m{1cm}|}
% Первая страница: полный заголовок
\caption{Параметры для настройки ЗП\hfill\vspace{-0.5\baselineskip}}\label{app_set:zp}\\
\hline
\rowcolor{gray!20}
Параметр & Условное обозначение на схеме & Значение/ Диапазон & Единица измерения & Шаг \\
\hline
\endfirsthead

% Промежуточные страницы: «Продолжение таблицы»
\caption*{\hspace{3pt}\emph{Продолжение таблицы \ref{app_set:zp}\hfill\vspace{-0.5\baselineskip}}} \\ % сделано по ГОСТ 2.105 п.6.8.7

\hline
\rowcolor{gray!20}
Параметр & Условное обозначение на схеме & Значение/ Диапазон & Единица измерения & Шаг \\
\hline
\endhead

%===m>ЗП|general+-
\multicolumn{5}{|c|}{ ЗП ВН } \\ \hline 
\raggedright Тср ЗП ВН & \centering 0 ... 300000 & \centering мс & \centering 5 & \centering \arraybackslash 300000 \\
\hline
\raggedright Iср ЗП ВН & \centering 0,04 ... 30,00 & \centering о.е. & \centering 0,01 & \centering \arraybackslash 30,00 \\
\hline
\raggedright ЗП ВН & \centering Введено \\ Выведено & \centering --- & \centering --- & \centering \arraybackslash Введено \\
\hline
\multicolumn{5}{|c|}{ ЗП ОО } \\ \hline 
\raggedright Тср ЗП ОО & \centering 0 ... 300000 & \centering мс & \centering 5 & \centering \arraybackslash 300000 \\
\hline
\raggedright Iср ЗП ОО & \centering 0,04 ... 30,00 & \centering о.е. & \centering 0,01 & \centering \arraybackslash 30,00 \\
\hline
\raggedright ЗП ОО & \centering Введено \\ Выведено & \centering --- & \centering --- & \centering \arraybackslash Введено \\
\hline
\multicolumn{5}{|c|}{ ЗП НН } \\ \hline 
\raggedright Тср ЗП НН & \centering 0 ... 300000 & \centering мс & \centering 5 & \centering \arraybackslash 300000 \\
\hline
\raggedright Iср ЗП НН & \centering 0,04 ... 30,00 & \centering о.е. & \centering 0,01 & \centering \arraybackslash 30,00 \\
\hline
\raggedright ЗП НН & \centering Введено \\ Выведено & \centering --- & \centering --- & \centering \arraybackslash Введено \\
\hline
%===m
\end{longtable}
\normalsize

%%%%%%%%%%%%%%%%%%%%%%%%%%%%%%%%%%%%%%% ЗПО

\par
Уставки ЗПО представлены в таблице \ref{app_set:zpo}.

\footnotesize
\begin{longtable}{|>{\centering\arraybackslash}m{5.3cm}|>{\centering\arraybackslash}m{3.3cm}|>{\centering\arraybackslash}m{4.2cm}|>{\centering\arraybackslash}m{1.8cm}|>{\centering\arraybackslash}m{1cm}|}
% Первая страница: полный заголовок
\caption{Параметры для настройки ЗПО\hfill\vspace{-0.5\baselineskip}}\label{app_set:zpo}\\
\hline
\rowcolor{gray!20}
Параметр & Условное обозначение на схеме & Значение/ Диапазон & Единица измерения & Шаг \\
\hline
\endfirsthead

% Промежуточные страницы: «Продолжение таблицы»
\caption*{\hspace{3pt}\emph{Продолжение таблицы \ref{app_set:zpo}\hfill\vspace{-0.5\baselineskip}}} \\ % сделано по ГОСТ 2.105 п.6.8.7

\hline
\rowcolor{gray!20}
Параметр & Условное обозначение на схеме & Значение/ Диапазон & Единица измерения & Шаг \\
\hline
\endhead

%===m>ЗПО АТ|general+-
\multicolumn{5}{|c|}{ ЗПО АТ } \\ \hline 
\raggedright Откл от ЗПО & \centering Введено \\ Выведено & \centering --- & \centering --- & \centering \arraybackslash Введено \\
\hline
\multicolumn{5}{|c|}{ ЗПО-1 } \\ \hline 
\raggedright ЗПО-1 конт ДТм & \centering Введено \\ Выведено & \centering --- & \centering --- & \centering \arraybackslash Введено \\
\hline
\raggedright Тср ЗПО-1 & \centering 0 ... 300000 & \centering мс & \centering 5 & \centering \arraybackslash 300000 \\
\hline
\raggedright ЗПО-1 контр РТ & \centering Выведено \\ От РТ1 \\ От РТ2 & \centering --- & \centering --- & \centering \arraybackslash Выведено \\
\hline
\raggedright ЗПО-1 & \centering Введено \\ Выведено & \centering --- & \centering --- & \centering \arraybackslash Введено \\
\hline
\multicolumn{5}{|c|}{ ЗПО-2 } \\ \hline 
\raggedright ЗПО-2 конт ДТм & \centering Введено \\ Выведено & \centering --- & \centering --- & \centering \arraybackslash Введено \\
\hline
\raggedright Тср ЗПО-2 & \centering 0 ... 300000 & \centering мс & \centering 5 & \centering \arraybackslash 300000 \\
\hline
\raggedright ЗПО-2 контр РТ & \centering Выведено \\ От РТ1 \\ От РТ2 & \centering --- & \centering --- & \centering \arraybackslash Выведено \\
\hline
\raggedright ЗПО-2 & \centering Введено \\ Выведено & \centering --- & \centering --- & \centering \arraybackslash Введено \\
\hline
\multicolumn{5}{|c|}{ ЗПО-3 } \\ \hline 
\raggedright ЗПО-3 конт ДТм & \centering Введено \\ Выведено & \centering --- & \centering --- & \centering \arraybackslash Введено \\
\hline
\raggedright Тср ЗПО-3 & \centering 0 ... 300000 & \centering мс & \centering 5 & \centering \arraybackslash 300000 \\
\hline
\raggedright ЗПО-3 контр РТ & \centering Выведено \\ От РТ1 \\ От РТ2 & \centering --- & \centering --- & \centering \arraybackslash Выведено \\
\hline
\raggedright ЗПО-3 & \centering Введено \\ Выведено & \centering --- & \centering --- & \centering \arraybackslash Введено \\
\hline
\multicolumn{5}{|c|}{ РТ ЗПО ВН } \\ \hline 
\raggedright Iср РТ1 ВН ЗПО & \centering 0,04 ... 30,00 & \centering о.е. & \centering 0,01 & \centering \arraybackslash 30,00 \\
\hline
\raggedright Iср РТ2 ВН ЗПО & \centering 0,04 ... 30,00 & \centering о.е. & \centering 0,01 & \centering \arraybackslash 30,00 \\
\hline
\raggedright РТ ВН ЗПО & \centering Введено \\ Выведено & \centering --- & \centering --- & \centering \arraybackslash Введено \\
\hline
\multicolumn{5}{|c|}{ РТ ЗПО ОО } \\ \hline 
\raggedright Iср РТ1 ОО ЗПО & \centering 0,04 ... 30,00 & \centering о.е. & \centering 0,01 & \centering \arraybackslash 30,00 \\
\hline
\raggedright Iср РТ2 ОО ЗПО & \centering 0,04 ... 30,00 & \centering о.е. & \centering 0,01 & \centering \arraybackslash 30,00 \\
\hline
\raggedright РТ ОО ЗПО & \centering Введено \\ Выведено & \centering --- & \centering --- & \centering \arraybackslash Введено \\
\hline
\multicolumn{5}{|c|}{ РТ ЗПО НН } \\ \hline 
\raggedright Iср РТ1 НН ЗПО & \centering 0,04 ... 30,00 & \centering о.е. & \centering 0,01 & \centering \arraybackslash 30,00 \\
\hline
\raggedright Iср РТ2 НН ЗПО & \centering 0,04 ... 30,00 & \centering о.е. & \centering 0,01 & \centering \arraybackslash 30,00 \\
\hline
\raggedright РТ НН ЗПО & \centering Введено \\ Выведено & \centering --- & \centering --- & \centering \arraybackslash Введено \\
\hline
%===m
\end{longtable}
\normalsize


%%%%%%%%%%%%%%%%%%%%%%%%%%%%%%%%%%%%%%% КИ НН

\par
Уставки КИ НН представлены в таблице \ref{app_set:kinn}.

\footnotesize
\begin{longtable}{|>{\centering\arraybackslash}m{5.3cm}|>{\centering\arraybackslash}m{3.3cm}|>{\centering\arraybackslash}m{4.2cm}|>{\centering\arraybackslash}m{1.8cm}|>{\centering\arraybackslash}m{1cm}|}
% Первая страница: полный заголовок
\caption{Параметры для настройки КИ НН\hfill\vspace{-0.5\baselineskip}}\label{app_set:kinn}\\
\hline
\rowcolor{gray!20}
Параметр & Условное обозначение на схеме & Значение/ Диапазон & Единица измерения & Шаг \\
\hline
\endfirsthead

% Промежуточные страницы: «Продолжение таблицы»
\caption*{\hspace{3pt}\emph{Продолжение таблицы \ref{app_set:kinn}\hfill\vspace{-0.5\baselineskip}}} \\ % сделано по ГОСТ 2.105 п.6.8.7

\hline
\rowcolor{gray!20}
Параметр & Условное обозначение на схеме & Значение/ Диапазон & Единица измерения & Шаг \\
\hline
\endhead

%===m>КИ НН|general+-
\multicolumn{5}{|c|}{ ИО КН } \\ \hline 
\raggedright 3U0ср КИ НН & \centering 0,0050 ... 1,5000 & \centering о.е. & \centering 0,0001 & \centering \arraybackslash 1,5000 \\
\hline
\raggedright Тср КИ НН & \centering 0 ... 300000 & \centering мс & \centering 5 & \centering \arraybackslash 300000 \\
\hline
\raggedright U2бл КИ НН & \centering 0,0050 ... 1,5000 & \centering о.е. & \centering 0,0001 & \centering \arraybackslash 1,5000 \\
\hline
\raggedright Блок U2 КИ НН & \centering Введено \\ Выведено & \centering --- & \centering --- & \centering \arraybackslash Введено \\
\hline
\raggedright КИ НН & \centering Введено \\ Выведено & \centering --- & \centering --- & \centering \arraybackslash Введено \\
\hline
%===m
\end{longtable}
\normalsize


%%%%%%%%%%%%%%%%%%%%%%%%%%%%%%%%%%%%%%% КЦН НН

\par
Уставки КЦН НН представлены в таблице \ref{app_set:kznnn1}.

\footnotesize
\begin{longtable}{|>{\centering\arraybackslash}m{5.3cm}|>{\centering\arraybackslash}m{3.3cm}|>{\centering\arraybackslash}m{4.2cm}|>{\centering\arraybackslash}m{1.8cm}|>{\centering\arraybackslash}m{1cm}|}
% Первая страница: полный заголовок
\caption{Параметры для настройки КЦН НН\hfill\vspace{-0.5\baselineskip}}\label{app_set:kznnn1}\\
\hline
\rowcolor{gray!20}
Параметр & Условное обозначение на схеме & Значение/ Диапазон & Единица измерения & Шаг \\
\hline
\endfirsthead

% Промежуточные страницы: «Продолжение таблицы»
\caption*{\hspace{3pt}\emph{Продолжение таблицы \ref{app_set:kznnn1}\hfill\vspace{-0.5\baselineskip}}} \\ % сделано по ГОСТ 2.105 п.6.8.7

\hline
\rowcolor{gray!20}
Параметр & Условное обозначение на схеме & Значение/ Диапазон & Единица измерения & Шаг \\
\hline
\endhead

%===m>КЦН|general+-
\multicolumn{5}{|c|}{ КЦН НН } \\ \hline 
\raggedright U2 КЦН НН & \centering 0,0050 ... 1,5000 & \centering о.е. & \centering 0,0001 & \centering \arraybackslash 1,5000 \\
\hline
\raggedright Uл КЦН НН & \centering 0,0050 ... 1,5000 & \centering о.е. & \centering 0,0001 & \centering \arraybackslash 0,0050 \\
\hline
\raggedright Тср КЦН НН & \centering 0 ... 300000 & \centering мс & \centering 5 & \centering \arraybackslash 300000 \\
\hline
\raggedright КЦН НН & \centering Введено \\ Выведено & \centering --- & \centering --- & \centering \arraybackslash Введено \\
\hline
\multicolumn{5}{|c|}{ КЦН НН1 } \\ \hline 
\raggedright U2 КЦН НН1 & \centering 0,0050 ... 1,5000 & \centering о.е. & \centering 0,0001 & \centering \arraybackslash 1,5000 \\
\hline
\raggedright Uл КЦН НН1 & \centering 0,0050 ... 1,5000 & \centering о.е. & \centering 0,0001 & \centering \arraybackslash 0,0050 \\
\hline
\raggedright Тср КЦН НН1 & \centering 0 ... 300000 & \centering мс & \centering 5 & \centering \arraybackslash 300000 \\
\hline
\raggedright  КЦН НН1 & \centering Введено \\ Выведено & \centering --- & \centering --- & \centering \arraybackslash Введено \\
\hline
\multicolumn{5}{|c|}{ КЦН НН2 } \\ \hline 
\raggedright U2 КЦН НН2 & \centering 0,0050 ... 1,5000 & \centering о.е. & \centering 0,0001 & \centering \arraybackslash 1,5000 \\
\hline
\raggedright Uл КЦН НН2 & \centering 0,0050 ... 1,5000 & \centering о.е. & \centering 0,0001 & \centering \arraybackslash 0,0050 \\
\hline
\raggedright Тср КЦН НН2 & \centering 0 ... 300000 & \centering мс & \centering 5 & \centering \arraybackslash 300000 \\
\hline
\raggedright  КЦН НН2 & \centering Введено \\ Выведено & \centering --- & \centering --- & \centering \arraybackslash Введено \\
\hline
%===m
\end{longtable}
\normalsize


%%%%%%%%%%%%%%%%%%%%%%%%%%%%%%%%%%%%%%% УРОВ ВН
\par
Уставки УРОВ ВН представлены в таблице \ref{app_set:urovvn}.

\footnotesize
\begin{longtable}{|>{\centering\arraybackslash}m{5.3cm}|>{\centering\arraybackslash}m{3.3cm}|>{\centering\arraybackslash}m{4.2cm}|>{\centering\arraybackslash}m{1.8cm}|>{\centering\arraybackslash}m{1cm}|}
% Первая страница: полный заголовок
\caption{Параметры для настройки УРОВ ВН\hfill\vspace{-0.5\baselineskip}}\label{app_set:urovvn}\\
\hline
\rowcolor{gray!20}
Параметр & Условное обозначение на схеме & Значение/ Диапазон & Единица измерения & Шаг \\
\hline
\endfirsthead

% Промежуточные страницы: «Продолжение таблицы»
\caption*{\hspace{3pt}\emph{Продолжение таблицы \ref{app_set:urovvn}\hfill\vspace{-0.5\baselineskip}}} \\ % сделано по ГОСТ 2.105 п.6.8.7

\hline
\rowcolor{gray!20}
Параметр & Условное обозначение на схеме & Значение/ Диапазон & Единица измерения & Шаг \\
\hline
\endhead

%===m>УРОВ В1(ЛВ) ВН|general+-
\multicolumn{5}{|c|}{ УРОВ } \\ \hline 
\raggedright Действ на себя УРОВ В1(ЛВ) ВН & \centering Введено \\ Выведено & \centering --- & \centering --- & \centering \arraybackslash Введено \\
\hline
\raggedright Тср повт УРОВ В1(ЛВ) ВН & \centering 0 ... 300000 & \centering мс & \centering 5 & \centering \arraybackslash 300000 \\
\hline
\raggedright Тср УРОВ В1(ЛВ) ВН & \centering 0 ... 300000 & \centering мс & \centering 5 & \centering \arraybackslash 300000 \\
\hline
\raggedright Iср УРОВ В1(ЛВ) ВН & \centering 0,04 ... 30,00 & \centering о.е. & \centering 0,01 & \centering \arraybackslash 30,00 \\
\hline
\raggedright УРОВ В1(ЛВ) ВН & \centering Введено \\ Выведено & \centering --- & \centering --- & \centering \arraybackslash Введено \\
\hline
\raggedright Подхв УРОВ В1(ЛВ) ВН & \centering Введено \\ Выведено & \centering --- & \centering --- & \centering \arraybackslash Введено \\
\hline
\raggedright Контр выкл УРОВ В1(ЛВ) ВН & \centering Введено \\ Выведено & \centering --- & \centering --- & \centering \arraybackslash Введено \\
\hline
%===m
\end{longtable}
\normalsize


%%%%%%%%%%%%%%%%%%%%%%%%%%%%%%%%%%%%%%% УРОВ СН
\par
Уставки УРОВ СН представлены в таблице \ref{app_set:urovsn}.

\footnotesize
\begin{longtable}{|>{\centering\arraybackslash}m{5.3cm}|>{\centering\arraybackslash}m{3.3cm}|>{\centering\arraybackslash}m{4.2cm}|>{\centering\arraybackslash}m{1.8cm}|>{\centering\arraybackslash}m{1cm}|}
% Первая страница: полный заголовок
\caption{Параметры для настройки УРОВ СН\hfill\vspace{-0.5\baselineskip}}\label{app_set:urovsn}\\
\hline
\rowcolor{gray!20}
Параметр & Условное обозначение на схеме & Значение/ Диапазон & Единица измерения & Шаг \\
\hline
\endfirsthead

% Промежуточные страницы: «Продолжение таблицы»
\caption*{\hspace{3pt}\emph{Продолжение таблицы \ref{app_set:urovsn}\hfill\vspace{-0.5\baselineskip}}} \\ % сделано по ГОСТ 2.105 п.6.8.7

\hline
\rowcolor{gray!20}
Параметр & Условное обозначение на схеме & Значение/ Диапазон & Единица измерения & Шаг \\
\hline
\endhead

%===m>УРОВ В1(ЛВ) СН|general+-
\multicolumn{5}{|c|}{ УРОВ } \\ \hline 
\raggedright Действ на себя УРОВ В1(ЛВ) СН & \centering Введено \\ Выведено & \centering --- & \centering --- & \centering \arraybackslash Введено \\
\hline
\raggedright Тср повт УРОВ В1(ЛВ) СН & \centering 0 ... 300000 & \centering мс & \centering 5 & \centering \arraybackslash 300000 \\
\hline
\raggedright Тср УРОВ В1(ЛВ) СН & \centering 0 ... 300000 & \centering мс & \centering 5 & \centering \arraybackslash 300000 \\
\hline
\raggedright Iср УРОВ В1(ЛВ) СН & \centering 0,04 ... 30,00 & \centering о.е. & \centering 0,01 & \centering \arraybackslash 30,00 \\
\hline
\raggedright УРОВ В1(ЛВ) СН & \centering Введено \\ Выведено & \centering --- & \centering --- & \centering \arraybackslash Введено \\
\hline
\raggedright Подхв УРОВ В1(ЛВ) СН & \centering Введено \\ Выведено & \centering --- & \centering --- & \centering \arraybackslash Введено \\
\hline
\raggedright Контр выкл УРОВ В1(ЛВ) СН & \centering Введено \\ Выведено & \centering --- & \centering --- & \centering \arraybackslash Введено \\
\hline
%===m
\end{longtable}
\normalsize



%%%%%%%%%%%%%%%%%%%%%%%%%%%%%%%%%%%%%%% ПО УРОВ НН
\par
Уставки ПО УРОВ НН представлены в таблице \ref{app_set:urovnn}.

\footnotesize
\begin{longtable}{|>{\centering\arraybackslash}m{5.3cm}|>{\centering\arraybackslash}m{3.3cm}|>{\centering\arraybackslash}m{4.2cm}|>{\centering\arraybackslash}m{1.8cm}|>{\centering\arraybackslash}m{1cm}|}
% Первая страница: полный заголовок
\caption{Параметры для настройки ПО УРОВ НН\hfill\vspace{-0.5\baselineskip}}\label{app_set:urovnn}\\
\hline
\rowcolor{gray!20}
Параметр & Условное обозначение на схеме & Значение/ Диапазон & Единица измерения & Шаг \\
\hline
\endfirsthead

% Промежуточные страницы: «Продолжение таблицы»
\caption*{\hspace{3pt}\emph{Продолжение таблицы \ref{app_set:urovnn}\hfill\vspace{-0.5\baselineskip}}} \\ % сделано по ГОСТ 2.105 п.6.8.7

\hline
\rowcolor{gray!20}
Параметр & Условное обозначение на схеме & Значение/ Диапазон & Единица измерения & Шаг \\
\hline
\endhead

%===m>ПО УРОВ НН|general+-
\multicolumn{5}{|c|}{ ПО УРОВ } \\ \hline 
\raggedright Тср ПО УРОВ НН & \centering 0 ... 300000 & \centering мс & \centering 5 & \centering \arraybackslash 300000 \\
\hline
\raggedright Iср ПО УРОВ НН & \centering 0,04 ... 30,00 & \centering о.е. & \centering 0,01 & \centering \arraybackslash 30,00 \\
\hline
\raggedright ПО УРОВ В НН & \centering Введено \\ Выведено & \centering --- & \centering --- & \centering \arraybackslash Выведено \\
\hline
\raggedright Подхв ПО УРОВ В НН & \centering Введено \\ Выведено & \centering --- & \centering --- & \centering \arraybackslash Выведено \\
\hline
%===m
\end{longtable}
\normalsize


%%%%%%%%%%%%%%%%%%%%%%%%%%%%%%%%%%%%%%% УВ ВН
\par
Уставки УВ ВН представлены в таблице \ref{app_set:uvvn}.

\footnotesize
\begin{longtable}{|>{\centering\arraybackslash}m{5.3cm}|>{\centering\arraybackslash}m{3.3cm}|>{\centering\arraybackslash}m{4.2cm}|>{\centering\arraybackslash}m{1.8cm}|>{\centering\arraybackslash}m{1cm}|}
% Первая страница: полный заголовок
\caption{Параметры для настройки УВ ВН\hfill\vspace{-0.5\baselineskip}}\label{app_set:uvvn}\\
\hline
\rowcolor{gray!20}
Параметр & Условное обозначение на схеме & Значение/ Диапазон & Единица измерения & Шаг \\
\hline
\endfirsthead

% Промежуточные страницы: «Продолжение таблицы»
\caption*{\hspace{3pt}\emph{Продолжение таблицы \ref{app_set:uvvn}\hfill\vspace{-0.5\baselineskip}}} \\ % сделано по ГОСТ 2.105 п.6.8.7

\hline
\rowcolor{gray!20}
Параметр & Условное обозначение на схеме & Значение/ Диапазон & Единица измерения & Шаг \\
\hline
\endhead

%===m>УВ В1(ЛВ) ВН|general+-
\multicolumn{5}{|c|}{ УВ } \\ \hline 
\raggedright Тожид усл опер вкл В1(ЛВ) ВН  & \centering 0 ... 300000 & \centering мс & \centering 5 & \centering \arraybackslash 10000 \\
\hline
\raggedright Тпрод опер вкл В1(ЛВ) ВН & \centering 0 ... 300000 & \centering мс & \centering 5 & \centering \arraybackslash 1000 \\
\hline
\raggedright Блок от многокр вкл на КЗ В1(ЛВ) ВН & \centering Введено \\ Выведено & \centering --- & \centering --- & \centering \arraybackslash Выведено \\
\hline
\raggedright Подхват вкл от РТ ЭМВ В1(ЛВ) ВН & \centering Введено \\ Выведено & \centering --- & \centering --- & \centering \arraybackslash Выведено \\
\hline
\raggedright УВ В1(ЛВ) ВН & \centering Введено \\ Выведено & \centering --- & \centering --- & \centering \arraybackslash Введено \\
\hline
%===m
\end{longtable}
\normalsize


%%%%%%%%%%%%%%%%%%%%%%%%%%%%%%%%%%%%%%% УВ СН
\par
Уставки УВ СН представлены в таблице \ref{app_set:uvsn}.

\footnotesize
\begin{longtable}{|>{\centering\arraybackslash}m{5.3cm}|>{\centering\arraybackslash}m{3.3cm}|>{\centering\arraybackslash}m{4.2cm}|>{\centering\arraybackslash}m{1.8cm}|>{\centering\arraybackslash}m{1cm}|}
% Первая страница: полный заголовок
\caption{Параметры для настройки УВ СН\hfill\vspace{-0.5\baselineskip}}\label{app_set:uvsn}\\
\hline
\rowcolor{gray!20}
Параметр & Условное обозначение на схеме & Значение/ Диапазон & Единица измерения & Шаг \\
\hline
\endfirsthead

% Промежуточные страницы: «Продолжение таблицы»
\caption*{\hspace{3pt}\emph{Продолжение таблицы \ref{app_set:uvsn}\hfill\vspace{-0.5\baselineskip}}} \\ % сделано по ГОСТ 2.105 п.6.8.7

\hline
\rowcolor{gray!20}
Параметр & Условное обозначение на схеме & Значение/ Диапазон & Единица измерения & Шаг \\
\hline
\endhead

%===m>УВ В1(ЛВ) СН|general+-
\multicolumn{5}{|c|}{ УВ } \\ \hline 
\raggedright Тожид усл опер вкл В1(ЛВ) СН & \centering 0 ... 300000 & \centering мс & \centering 5 & \centering \arraybackslash 10000 \\
\hline
\raggedright Тпрод опер вкл В1(ЛВ) СН & \centering 0 ... 300000 & \centering мс & \centering 5 & \centering \arraybackslash 1000 \\
\hline
\raggedright Блок от многокр вкл на КЗ В1(ЛВ) СН & \centering Введено \\ Выведено & \centering --- & \centering --- & \centering \arraybackslash Выведено \\
\hline
\raggedright Подхват вкл от РТ ЭМВ В1(ЛВ) СН & \centering Введено \\ Выведено & \centering --- & \centering --- & \centering \arraybackslash Выведено \\
\hline
\raggedright УВ В1(ЛВ) СН & \centering Введено \\ Выведено & \centering --- & \centering --- & \centering \arraybackslash Введено \\
\hline
%===m
\end{longtable}
\normalsize


%%%%%%%%%%%%%%%%%%%%%%%%%%%%%%%%%%%%%%% КСН ВН
\par
Уставки КСН ВН представлены в таблице \ref{app_set:ksnvn}.

\footnotesize
\begin{longtable}{|>{\centering\arraybackslash}m{5.3cm}|>{\centering\arraybackslash}m{3.3cm}|>{\centering\arraybackslash}m{4.2cm}|>{\centering\arraybackslash}m{1.8cm}|>{\centering\arraybackslash}m{1cm}|}
% Первая страница: полный заголовок
\caption{Параметры для настройки КСН ВН\hfill\vspace{-0.5\baselineskip}}\label{app_set:ksnvn}\\
\hline
\rowcolor{gray!20}
Параметр & Условное обозначение на схеме & Значение/ Диапазон & Единица измерения & Шаг \\
\hline
\endfirsthead

% Промежуточные страницы: «Продолжение таблицы»
\caption*{\hspace{3pt}\emph{Продолжение таблицы \ref{app_set:ksnvn}\hfill\vspace{-0.5\baselineskip}}} \\ % сделано по ГОСТ 2.105 п.6.8.7

\hline
\rowcolor{gray!20}
Параметр & Условное обозначение на схеме & Значение/ Диапазон & Единица измерения & Шаг \\
\hline
\endhead

%===m>КСН В1(ЛВ) ВН|general+-
\multicolumn{5}{|c|}{ КСН } \\ \hline 
\raggedright Uотс Т & \centering 0,0050 ... 1,5000 & \centering о.е. & \centering 0,0001 & \centering \arraybackslash 0,2000 \\
\hline
\raggedright Uнал СШ & \centering 0,0050 ... 1,5000 & \centering о.е. & \centering 0,0001 & \centering \arraybackslash 0,8000 \\
\hline
\raggedright Uотс СШ & \centering 0,0050 ... 1,5000 & \centering о.е. & \centering 0,0001 & \centering \arraybackslash 0,2000 \\
\hline
\raggedright Uнал Т & \centering 0,0050 ... 1,5000 & \centering о.е. & \centering 0,0001 & \centering \arraybackslash 0,8000 \\
\hline
\raggedright КС & \centering Введено \\ Выведено & \centering --- & \centering --- & \centering \arraybackslash Введено \\
\hline
\raggedright U<Uуст & \centering 0,0050 ... 1,5000 & \centering о.е. & \centering 0,0001 & \centering \arraybackslash 0,1000 \\
\hline
\raggedright df макс & \centering 0,050 ... 3,000 & \centering Гц & \centering 0,001 & \centering \arraybackslash 0,400 \\
\hline
\raggedright dФ макс & \centering 1 ... 90 & \centering Эл. градусы & \centering 1,0 & \centering \arraybackslash 10 \\
\hline
\raggedright УС & \centering Введено \\ Выведено & \centering --- & \centering --- & \centering \arraybackslash Введено \\
\hline
\raggedright df доп & \centering 0,025 ... 3,000 & \centering Гц & \centering 0,001 & \centering \arraybackslash 1,000 \\
\hline
\raggedright Т опер. вкл & \centering 0 ... 10000 & \centering мс & \centering 5 & \centering \arraybackslash 100 \\
\hline
\raggedright Контр опер вкл & \centering Внутр \\ Внеш & \centering --- & \centering --- & \centering \arraybackslash Внутр \\
\hline
\raggedright Опер вкл с КОН & \centering Введено \\ Выведено & \centering --- & \centering --- & \centering \arraybackslash Введено \\
\hline
\raggedright Опер вкл В без контролей & \centering Введено \\ Выведено & \centering --- & \centering --- & \centering \arraybackslash Введено \\
\hline
\raggedright КСН & \centering Введено \\ Выведено & \centering --- & \centering --- & \centering \arraybackslash Введено \\
\hline
%===m
\end{longtable}
\normalsize



%%%%%%%%%%%%%%%%%%%%%%%%%%%%%%%%%%%%%%% КСН СН
\par
Уставки КСН СН представлены в таблице \ref{app_set:ksnsn}.

\footnotesize
\begin{longtable}{|>{\centering\arraybackslash}m{5.3cm}|>{\centering\arraybackslash}m{3.3cm}|>{\centering\arraybackslash}m{4.2cm}|>{\centering\arraybackslash}m{1.8cm}|>{\centering\arraybackslash}m{1cm}|}
% Первая страница: полный заголовок
\caption{Параметры для настройки КСН СН\hfill\vspace{-0.5\baselineskip}}\label{app_set:ksnsn}\\
\hline
\rowcolor{gray!20}
Параметр & Условное обозначение на схеме & Значение/ Диапазон & Единица измерения & Шаг \\
\hline
\endfirsthead

% Промежуточные страницы: «Продолжение таблицы»
\caption*{\hspace{3pt}\emph{Продолжение таблицы \ref{app_set:ksnsn}\hfill\vspace{-0.5\baselineskip}}} \\ % сделано по ГОСТ 2.105 п.6.8.7

\hline
\rowcolor{gray!20}
Параметр & Условное обозначение на схеме & Значение/ Диапазон & Единица измерения & Шаг \\
\hline
\endhead

%===m>КСН В1(ЛВ) СН|general+-
\multicolumn{5}{|c|}{ КСН } \\ \hline 
\raggedright Uотс Т & \centering 0,0050 ... 1,5000 & \centering о.е. & \centering 0,0001 & \centering \arraybackslash 0,2000 \\
\hline
\raggedright Uнал СШ & \centering 0,0050 ... 1,5000 & \centering о.е. & \centering 0,0001 & \centering \arraybackslash 0,8000 \\
\hline
\raggedright Uотс СШ & \centering 0,0050 ... 1,5000 & \centering о.е. & \centering 0,0001 & \centering \arraybackslash 0,2000 \\
\hline
\raggedright Uнал Т & \centering 0,0050 ... 1,5000 & \centering о.е. & \centering 0,0001 & \centering \arraybackslash 0,8000 \\
\hline
\raggedright КС & \centering Введено \\ Выведено & \centering --- & \centering --- & \centering \arraybackslash Введено \\
\hline
\raggedright U<Uуст & \centering 0,0050 ... 1,5000 & \centering о.е. & \centering 0,0001 & \centering \arraybackslash 0,1000 \\
\hline
\raggedright df макс & \centering 0,050 ... 3,000 & \centering Гц & \centering 0,001 & \centering \arraybackslash 0,400 \\
\hline
\raggedright dФ макс & \centering 1 ... 90 & \centering Эл. градусы & \centering 1,0 & \centering \arraybackslash 10 \\
\hline
\raggedright УС & \centering Введено \\ Выведено & \centering --- & \centering --- & \centering \arraybackslash Введено \\
\hline
\raggedright df доп & \centering 0,025 ... 3,000 & \centering Гц & \centering 0,001 & \centering \arraybackslash 1,000 \\
\hline
\raggedright Т опер. вкл & \centering 0 ... 10000 & \centering мс & \centering 5 & \centering \arraybackslash 100 \\
\hline
\raggedright Контр опер вкл & \centering Внутр \\ Внеш & \centering --- & \centering --- & \centering \arraybackslash Внутр \\
\hline
\raggedright Опер вкл с КОН & \centering Введено \\ Выведено & \centering --- & \centering --- & \centering \arraybackslash Введено \\
\hline
\raggedright Опер вкл В без контролей & \centering Введено \\ Выведено & \centering --- & \centering --- & \centering \arraybackslash Введено \\
\hline
\raggedright КСН & \centering Введено \\ Выведено & \centering --- & \centering --- & \centering \arraybackslash Введено \\
\hline
%===m
\end{longtable}
\normalsize


%%%%%%%%%%%%%%%%%%%%%%%%%%%%%%%%%%%%%%% АПВ ВН
\par
Уставки АПВ ВН представлены в таблице \ref{app_set:apvvn}.

\footnotesize
\begin{longtable}{|>{\centering\arraybackslash}m{5.3cm}|>{\centering\arraybackslash}m{3.3cm}|>{\centering\arraybackslash}m{4.2cm}|>{\centering\arraybackslash}m{1.8cm}|>{\centering\arraybackslash}m{1cm}|}
% Первая страница: полный заголовок
\caption{Параметры для настройки АПВ ВН\hfill\vspace{-0.5\baselineskip}}\label{app_set:apvvn}\\
\hline
\rowcolor{gray!20}
Параметр & Условное обозначение на схеме & Значение/ Диапазон & Единица измерения & Шаг \\
\hline
\endfirsthead

% Промежуточные страницы: «Продолжение таблицы»
\caption*{\hspace{3pt}\emph{Продолжение таблицы \ref{app_set:apvvn}\hfill\vspace{-0.5\baselineskip}}} \\ % сделано по ГОСТ 2.105 п.6.8.7

\hline
\rowcolor{gray!20}
Параметр & Условное обозначение на схеме & Значение/ Диапазон & Единица измерения & Шаг \\
\hline
\endhead

%===m>АПВ В1(ЛВ) ВН|general+-
\multicolumn{5}{|c|}{ АПВ В1(ЛВ) ВН } \\ \hline 
\raggedright Тгот АПВ В1(ЛВ) ВН & \centering 0 ... 300000 & \centering мс & \centering 5 & \centering \arraybackslash 10000 \\
\hline
\raggedright Тожид усл вкл В1(ЛВ) ВН & \centering 0 ... 300000 & \centering мс & \centering 5 & \centering \arraybackslash 20000 \\
\hline
\multicolumn{5}{|c|}{ АПВ СШ } \\ \hline 
\raggedright АПВ СШ В1(ЛВ) ВН & \centering Введено \\ Выведено & \centering --- & \centering --- & \centering \arraybackslash Выведено \\
\hline
\raggedright Тгот АПВ & \centering 0 ... 300000 & \centering мс & \centering 5 & \centering \arraybackslash 15000 \\
\hline
\raggedright Тожид усл. вкл. & \centering 0 ... 300000 & \centering мс & \centering 5 & \centering \arraybackslash 100 \\
\hline
\raggedright Тср АПВш В1(ЛВ) ВН & \centering 0 ... 300000 & \centering мс & \centering 5 & \centering \arraybackslash 2000 \\
\hline
\raggedright Тимп АПВш В1(ЛВ) ВН & \centering 0 ... 300000 & \centering мс & \centering 5 & \centering \arraybackslash 20 \\
\hline
\multicolumn{5}{|c|}{ АПВ Т } \\ \hline 
\raggedright Тгот АПВ & \centering 0 ... 300000 & \centering мс & \centering 5 & \centering \arraybackslash 10000 \\
\hline
\raggedright Тожид усл вкл & \centering 0 ... 300000 & \centering мс & \centering 5 & \centering \arraybackslash 20000 \\
\hline
\raggedright АПВ В1(ЛВ) ВН & \centering Введено \\ Выведено & \centering --- & \centering --- & \centering \arraybackslash Выведено \\
\hline
\raggedright Тср АПВ-1 В1(ЛВ) ВН & \centering 0 ... 300000 & \centering мс & \centering 5 & \centering \arraybackslash 1000 \\
\hline
\raggedright Тимп АПВ В1(ЛВ) ВН & \centering 0 ... 300000 & \centering мс & \centering 5 & \centering \arraybackslash 2000 \\
\hline
\raggedright Тср АПВ-2 В1(ЛВ) ВН & \centering 0 ... 300000 & \centering мс & \centering 5 & \centering \arraybackslash 4000 \\
\hline
\raggedright Твосст гот АПВ В1(ЛВ) ВН & \centering 0 ... 300000 & \centering мс & \centering 5 & \centering \arraybackslash 10000 \\
\hline
\raggedright 2-ой цикл АПВ В1(ЛВ) ВН  & \centering Введено \\ Выведено & \centering --- & \centering --- & \centering \arraybackslash Выведено \\
\hline
%===m
\end{longtable}
\normalsize


%%%%%%%%%%%%%%%%%%%%%%%%%%%%%%%%%%%%%%% АПВ СН
\par
Уставки АПВ СН представлены в таблице \ref{app_set:apvsn}.

\footnotesize
\begin{longtable}{|>{\centering\arraybackslash}m{5.3cm}|>{\centering\arraybackslash}m{3.3cm}|>{\centering\arraybackslash}m{4.2cm}|>{\centering\arraybackslash}m{1.8cm}|>{\centering\arraybackslash}m{1cm}|}
% Первая страница: полный заголовок
\caption{Параметры для настройки АПВ СН\hfill\vspace{-0.5\baselineskip}}\label{app_set:apvsn}\\
\hline
\rowcolor{gray!20}
Параметр & Условное обозначение на схеме & Значение/ Диапазон & Единица измерения & Шаг \\
\hline
\endfirsthead

% Промежуточные страницы: «Продолжение таблицы»
\caption*{\hspace{3pt}\emph{Продолжение таблицы \ref{app_set:apvsn}\hfill\vspace{-0.5\baselineskip}}} \\ % сделано по ГОСТ 2.105 п.6.8.7

\hline
\rowcolor{gray!20}
Параметр & Условное обозначение на схеме & Значение/ Диапазон & Единица измерения & Шаг \\
\hline
\endhead

%===m>АПВ В1(ЛВ) СН|general+-
\multicolumn{5}{|c|}{ АПВ В1(ЛВ) СН } \\ \hline 
\raggedright Тгот АПВ В1(ЛВ) СН & \centering 0 ... 300000 & \centering мс & \centering 5 & \centering \arraybackslash 10000 \\
\hline
\raggedright Тожид усл вкл В1(ЛВ) СН & \centering 0 ... 300000 & \centering мс & \centering 5 & \centering \arraybackslash 20000 \\
\hline
\multicolumn{5}{|c|}{ АПВ СШ } \\ \hline 
\raggedright АПВ СШ В1(ЛВ) СН & \centering Введено \\ Выведено & \centering --- & \centering --- & \centering \arraybackslash Выведено \\
\hline
\raggedright Тгот АПВ & \centering 0 ... 300000 & \centering мс & \centering 5 & \centering \arraybackslash 15000 \\
\hline
\raggedright Тожид усл. вкл. & \centering 0 ... 300000 & \centering мс & \centering 5 & \centering \arraybackslash 100 \\
\hline
\raggedright Тср АПВш В1(ЛВ) СН & \centering 0 ... 300000 & \centering мс & \centering 5 & \centering \arraybackslash 2000 \\
\hline
\raggedright Тимп АПВш В1(ЛВ) СН & \centering 0 ... 300000 & \centering мс & \centering 5 & \centering \arraybackslash 20 \\
\hline
\multicolumn{5}{|c|}{ АПВ Т } \\ \hline 
\raggedright Тгот АПВ & \centering 0 ... 300000 & \centering мс & \centering 5 & \centering \arraybackslash 10000 \\
\hline
\raggedright Тожид усл вкл & \centering 0 ... 300000 & \centering мс & \centering 5 & \centering \arraybackslash 20000 \\
\hline
\raggedright АПВ В1(ЛВ) СН & \centering Введено \\ Выведено & \centering --- & \centering --- & \centering \arraybackslash Выведено \\
\hline
\raggedright Тср АПВ-1 В1(ЛВ) СН & \centering 0 ... 300000 & \centering мс & \centering 5 & \centering \arraybackslash 1000 \\
\hline
\raggedright Тимп АПВ В1(ЛВ) СН & \centering 0 ... 300000 & \centering мс & \centering 5 & \centering \arraybackslash 2000 \\
\hline
\raggedright Тср АПВ-2 В1(ЛВ) СН & \centering 0 ... 300000 & \centering мс & \centering 5 & \centering \arraybackslash 4000 \\
\hline
\raggedright Твосст гот АПВ В1(ЛВ) СН & \centering 0 ... 300000 & \centering мс & \centering 5 & \centering \arraybackslash 10000 \\
\hline
\raggedright 2-ой цикл АПВ В1(ЛВ) СН & \centering Введено \\ Выведено & \centering --- & \centering --- & \centering \arraybackslash Выведено \\
\hline
%===m
\end{longtable}
\normalsize




%%%%%%%%%%%%%%%%%%%%%%%%%%%%%%%%%%%%%%%%%%%%%%%%%%%%%%%%%%%%%%%%%%%%%%%%% КСВ1 %%%%%%%%%%%%%%%%%%%%%%%%%%%%%%%%%%%%%%%%%%%%%%%%%%%%%%%%%%%%%%%%%%%%%%%%%%%

\par
Уставки КСВ В1(ЛВ) ВН представлены в таблице \ref{app_set:ksv1}.

\footnotesize
\begin{longtable}{|>{\centering\arraybackslash}m{5.3cm}|>{\centering\arraybackslash}m{3.3cm}|>{\centering\arraybackslash}m{4.2cm}|>{\centering\arraybackslash}m{1.8cm}|>{\centering\arraybackslash}m{1cm}|}
\caption{Параметры для настройки КСВ В1\hfill\vspace{-0.5\baselineskip}}\label{app_set:ksv1}\\ 
\hline
\rowcolor{gray!20}
Параметр & Условное обозначение на схеме & Значение/ Диапазон & Единица измерения & Шаг \\ 
\hline
\endfirsthead
\caption*{\hspace{3pt}\emph{Продолжение таблицы \ref{app_set:ksv1}\hfill\vspace{-0.5\baselineskip}}} \\ % сделано по ГОСТ 2.105 п.6.8.7
\hline
\rowcolor{gray!20}
Параметр & Условное обозначение на схеме & Значение/ Диапазон & Единица измерения & Шаг \\ 
%\hline
\endhead

%===m>КСВ В1(ЛВ) ВН|general+-
\multicolumn{5}{|c|}{ КСВ В1(ЛВ) ВН } \\ \hline 
\raggedright Кол-во ЭМО В1(ЛВ) ВН & \centering Один \\ Два & \centering --- & \centering --- & \centering \arraybackslash Два \\
\hline
\raggedright КСВ В1(ЛВ) ВН & \centering Введено \\ Выведено & \centering --- & \centering --- & \centering \arraybackslash Введено \\
\hline
\multicolumn{5}{|c|}{ Контроль В } \\ \hline 
\raggedright Контр ЭМО в любом пол В1(ЛВ) ВН & \centering Выведено \\ ЭМО1 \\ ЭМО2 \\ ЭМО1 и ЭМО2 & \centering --- & \centering --- & \centering \arraybackslash Выведено \\
\hline
\raggedright Кол-во ЭМО В & \centering Один \\ Два & \centering --- & \centering --- & \centering \arraybackslash Два \\
\hline
\raggedright Контр ОТ сигн В1(ЛВ) ВН & \centering Введено \\ Выведено & \centering --- & \centering --- & \centering \arraybackslash Введено \\
\hline
\raggedright Тср Неиспр ЦУ В1(ЛВ) ВН & \centering 0 ... 300000 & \centering мс & \centering 5 & \centering \arraybackslash 1000 \\
\hline
\raggedright Тср Пруж не заведены В1(ЛВ) ВН & \centering 0 ... 300000 & \centering мс & \centering 5 & \centering \arraybackslash 1000 \\
\hline
\raggedright Тср Неиспр обогр В1(ЛВ) ВН & \centering 0 ... 300000 & \centering мс & \centering 5 & \centering \arraybackslash 1000 \\
\hline
\multicolumn{5}{|c|}{ Контроль ТТ } \\ \hline 
\raggedright Откл от Авар. давл. SF6 ТТ В1(ЛВ) ВН & \centering Выведено \\ Введено & \centering --- & \centering --- & \centering \arraybackslash Введено \\
\hline
\raggedright Тср Неиспр. обогр. ТТ В1(ЛВ) ВН & \centering 0 ... 300000 & \centering мс & \centering 5 & \centering \arraybackslash 15000 \\
\hline
\multicolumn{5}{|c|}{ РФ } \\ \hline 
\raggedright Реле фиксации В1(ЛВ) ВН & \centering РФК \\ РФП & \centering --- & \centering --- & \centering \arraybackslash РФК \\
\hline
\multicolumn{5}{|c|}{ Защита ЭМУ } \\ \hline 
\raggedright Обест. ЭМУ от Авар. давл. SF6 В1(ЛВ) ВН & \centering Введено \\ Выведено & \centering --- & \centering --- & \centering \arraybackslash Введено \\
\hline
\raggedright Тср защиты ЭМВ В1(ЛВ) ВН & \centering 0 ... 300000 & \centering мс & \centering 5 & \centering \arraybackslash 1000 \\
\hline
\raggedright Тср защиты ЭМО1 В1(ЛВ) ВН & \centering 0 ... 300000 & \centering мс & \centering 5 & \centering \arraybackslash 1000 \\
\hline
\raggedright Тср защиты ЭМО2 В1(ЛВ) ВН & \centering 0 ... 300000 & \centering мс & \centering 5 & \centering \arraybackslash 1000 \\
\hline
%===m
\end{longtable}
\normalsize

%%%%%%%%%%%%%%%%%%%%%%%%%%%%%%%%%%%%%%%%%%%%%%%%%%%%%%%%%%%%%%%%%%%%%%%%% КСВ2 %%%%%%%%%%%%%%%%%%%%%%%%%%%%%%%%%%%%%%%%%%%%%%%%%%%%%%%%%%%%%%%%%%%%%%%%%%%

\par
Уставки КСВ В1(ЛВ) СН представлены в таблице \ref{app_set:ksv2}.

\footnotesize
\begin{longtable}{|>{\centering\arraybackslash}m{5.3cm}|>{\centering\arraybackslash}m{3.3cm}|>{\centering\arraybackslash}m{4.2cm}|>{\centering\arraybackslash}m{1.8cm}|>{\centering\arraybackslash}m{1cm}|}
\caption{Параметры для настройки КСВ В1\hfill\vspace{-0.5\baselineskip}}\label{app_set:ksv2}\\ 
\hline
\rowcolor{gray!20}
Параметр & Условное обозначение на схеме & Значение/ Диапазон & Единица измерения & Шаг \\ 
\hline
\endfirsthead
\caption*{\hspace{3pt}\emph{Продолжение таблицы \ref{app_set:ksv2}\hfill\vspace{-0.5\baselineskip}}} \\ % сделано по ГОСТ 2.105 п.6.8.7
\hline
\rowcolor{gray!20}
Параметр & Условное обозначение на схеме & Значение/ Диапазон & Единица измерения & Шаг \\ 
%\hline
\endhead

%===m>КСВ В1(ЛВ) СН|general+-
\multicolumn{5}{|c|}{ КСВ В1(ЛВ) СН } \\ \hline 
\raggedright Кол-во ЭМО В1(ЛВ) СН & \centering Один \\ Два & \centering --- & \centering --- & \centering \arraybackslash Два \\
\hline
\raggedright КСВ В1(ЛВ) СН & \centering Введено \\ Выведено & \centering --- & \centering --- & \centering \arraybackslash Введено \\
\hline
\multicolumn{5}{|c|}{ Контроль В } \\ \hline 
\raggedright Контр. ЭМО в любом пол. В1(ЛВ) СН & \centering Выведено \\ ЭМО1 \\ ЭМО2 \\ ЭМО1 и ЭМО2 & \centering --- & \centering --- & \centering \arraybackslash Выведено \\
\hline
\raggedright Кол-во ЭМО В & \centering Один \\ Два & \centering --- & \centering --- & \centering \arraybackslash Два \\
\hline
\raggedright Контр. ОТ сигн. В1(ЛВ) СН & \centering Введено \\ Выведено & \centering --- & \centering --- & \centering \arraybackslash Введено \\
\hline
\raggedright Тср Неиспр. ЦУ В1(ЛВ) СН & \centering 0 ... 300000 & \centering мс & \centering 5 & \centering \arraybackslash 1000 \\
\hline
\raggedright Тср Пруж. не заведены В1(ЛВ) СН & \centering 0 ... 300000 & \centering мс & \centering 5 & \centering \arraybackslash 1000 \\
\hline
\raggedright Тср Неиспр. обогр. В1(ЛВ) СН & \centering 0 ... 300000 & \centering мс & \centering 5 & \centering \arraybackslash 1000 \\
\hline
\multicolumn{5}{|c|}{ Контроль ТТ } \\ \hline 
\raggedright Откл от Авар. давл. SF6 ТТ В1(ЛВ) СН & \centering Выведено \\ Введено & \centering --- & \centering --- & \centering \arraybackslash Выведено \\
\hline
\raggedright Тср Неиспр. обогр. ТТ В1(ЛВ) СН & \centering 0 ... 300000 & \centering мс & \centering 5 & \centering \arraybackslash 15000 \\
\hline
\multicolumn{5}{|c|}{ РФ } \\ \hline 
\raggedright Реле фиксации В1(ЛВ) СН & \centering РФК \\ РФП & \centering --- & \centering --- & \centering \arraybackslash РФК \\
\hline
\multicolumn{5}{|c|}{ Защита ЭМУ } \\ \hline 
\raggedright Обест. ЭМУ от Авар. давл. SF6 В1(ЛВ) СН & \centering Введено \\ Выведено & \centering --- & \centering --- & \centering \arraybackslash Введено \\
\hline
\raggedright Тср защиты ЭМВ В1(ЛВ) СН & \centering 0 ... 300000 & \centering мс & \centering 5 & \centering \arraybackslash 1000 \\
\hline
\raggedright Тср защиты ЭМО1 В1(ЛВ) СН & \centering 0 ... 300000 & \centering мс & \centering 5 & \centering \arraybackslash 1000 \\
\hline
\raggedright Тср защиты ЭМО2 В1(ЛВ) СН & \centering 0 ... 300000 & \centering мс & \centering 5 & \centering \arraybackslash 1000 \\
\hline
%===m
\end{longtable}
\normalsize



%%%%%%%%%%%%%%%%%%%%%%%%%%%%%%%%%%%%%%%%%%%%%%%%%%%%%%%%%%%%%%%%%%%%%%%%% ЛО внеш РЗ ВН %%%%%%%%%%%%%%%%%%%%%%%%%%%%%%%%%%%%%%%%%%%%%%%%%%%%%%%%%%%%%%%%%%%%%%%%%%%

\par
Уставки ЛО внеш РЗ ВН представлены в таблице \ref{app_set:lovneshvn}.

\footnotesize
\begin{longtable}{|>{\centering\arraybackslash}m{5.3cm}|>{\centering\arraybackslash}m{3.3cm}|>{\centering\arraybackslash}m{4.2cm}|>{\centering\arraybackslash}m{1.8cm}|>{\centering\arraybackslash}m{1cm}|}
\caption{Параметры для настройки ЛО внеш РЗ ВН\hfill\vspace{-0.5\baselineskip}}\label{app_set:lovneshvn}\\ 
\hline
\rowcolor{gray!20}
Параметр & Условное обозначение на схеме & Значение/ Диапазон & Единица измерения & Шаг \\ 
\hline
\endfirsthead
\caption*{\hspace{3pt}\emph{Продолжение таблицы \ref{app_set:lovneshvn}\hfill\vspace{-0.5\baselineskip}}} \\ % сделано по ГОСТ 2.105 п.6.8.7
\hline
\rowcolor{gray!20}
Параметр & Условное обозначение на схеме & Значение/ Диапазон & Единица измерения & Шаг \\ 
%\hline
\endhead

%===m>ЛО внеш РЗ ВН|general+-
\multicolumn{5}{|c|}{ ЛО внеш РЗ ВН(СН) } \\ \hline 
\raggedright ЛО внеш РЗ ВН & \centering Введено \\ Выведено & \centering --- & \centering --- & \centering \arraybackslash Введено \\
\hline
%===m
\end{longtable}
\normalsize


%%%%%%%%%%%%%%%%%%%%%%%%%%%%%%%%%%%%%%%%%%%%%%%%%%%%%%%%%%%%%%%%%%%%%%%%% ЛО внеш РЗ СН %%%%%%%%%%%%%%%%%%%%%%%%%%%%%%%%%%%%%%%%%%%%%%%%%%%%%%%%%%%%%%%%%%%%%%%%%%%

\par
Уставки ЛО внеш РЗ СН представлены в таблице \ref{app_set:lovneshsn}.

\footnotesize
\begin{longtable}{|>{\centering\arraybackslash}m{5.3cm}|>{\centering\arraybackslash}m{3.3cm}|>{\centering\arraybackslash}m{4.2cm}|>{\centering\arraybackslash}m{1.8cm}|>{\centering\arraybackslash}m{1cm}|}
\caption{Параметры для настройки ЛО внеш РЗ СН\hfill\vspace{-0.5\baselineskip}}\label{app_set:lovneshsn}\\ 
\hline
\rowcolor{gray!20}
Параметр & Условное обозначение на схеме & Значение/ Диапазон & Единица измерения & Шаг \\ 
\hline
\endfirsthead
\caption*{\hspace{3pt}\emph{Продолжение таблицы \ref{app_set:lovneshsn}\hfill\vspace{-0.5\baselineskip}}} \\ % сделано по ГОСТ 2.105 п.6.8.7
\hline
\rowcolor{gray!20}
Параметр & Условное обозначение на схеме & Значение/ Диапазон & Единица измерения & Шаг \\ 
%\hline
\endhead

%===m>ЛО внеш РЗ СН|general+-
\multicolumn{5}{|c|}{ ЛО внеш РЗ ВН(СН) } \\ \hline 
\raggedright ЛО внеш РЗ СН & \centering Введено \\ Выведено & \centering --- & \centering --- & \centering \arraybackslash Введено \\
\hline
%===m
\end{longtable}
\normalsize


%%%%%%%%%%%%%%%%%%%%%%%%%%%%%%%%%%%%%%%%%%%%%%%%%%%%%%%%%%%%%%%%%%%%%%%%% ЛО внеш РЗ НН %%%%%%%%%%%%%%%%%%%%%%%%%%%%%%%%%%%%%%%%%%%%%%%%%%%%%%%%%%%%%%%%%%%%%%%%%%% "ЛО внеш АПТ"

\par
Уставки ЛО внеш РЗ НН представлены в таблице \ref{app_set:lovneshnn}.

\footnotesize
\begin{longtable}{|>{\centering\arraybackslash}m{5.3cm}|>{\centering\arraybackslash}m{3.3cm}|>{\centering\arraybackslash}m{4.2cm}|>{\centering\arraybackslash}m{1.8cm}|>{\centering\arraybackslash}m{1cm}|}
\caption{Параметры для настройки ЛО внеш РЗ НН\hfill\vspace{-0.5\baselineskip}}\label{app_set:lovneshnn}\\ 
\hline
\rowcolor{gray!20}
Параметр & Условное обозначение на схеме & Значение/ Диапазон & Единица измерения & Шаг \\ 
\hline
\endfirsthead
\caption*{\hspace{3pt}\emph{Продолжение таблицы \ref{app_set:lovneshnn}\hfill\vspace{-0.5\baselineskip}}} \\ % сделано по ГОСТ 2.105 п.6.8.7
\hline
\rowcolor{gray!20}
Параметр & Условное обозначение на схеме & Значение/ Диапазон & Единица измерения & Шаг \\ 
%\hline
\endhead

%===m>ЛО внеш РЗ НН|general+-
\multicolumn{5}{|c|}{ ЛО внеш РЗ НН } \\ \hline 
\raggedright ЛО внеш РЗ НН & \centering Введено \\ Выведено & \centering --- & \centering --- & \centering \arraybackslash Введено \\
\hline
%===m
\end{longtable}
\normalsize



%%%%%%%%%%%%%%%%%%%%%%%%%%%%%%%%%%%%%%%%%%%%%%%%%%%%%%%%%%%%%%%%%%%%%%%%% ЛО внеш АПТ %%%%%%%%%%%%%%%%%%%%%%%%%%%%%%%%%%%%%%%%%%%%%%%%%%%%%%%%%%%%%%%%%%%%%%%%%%%

\par
Уставки ЛО внеш АПТ представлены в таблице \ref{app_set:lovneshapt}.

\footnotesize
\begin{longtable}{|>{\centering\arraybackslash}m{5.3cm}|>{\centering\arraybackslash}m{3.3cm}|>{\centering\arraybackslash}m{4.2cm}|>{\centering\arraybackslash}m{1.8cm}|>{\centering\arraybackslash}m{1cm}|}
\caption{Параметры для настройки ЛО внеш АПТ\hfill\vspace{-0.5\baselineskip}}\label{app_set:lovneshapt}\\ 
\hline
\rowcolor{gray!20}
Параметр & Условное обозначение на схеме & Значение/ Диапазон & Единица измерения & Шаг \\ 
\hline
\endfirsthead
\caption*{\hspace{3pt}\emph{Продолжение таблицы \ref{app_set:lovneshapt}\hfill\vspace{-0.5\baselineskip}}} \\ % сделано по ГОСТ 2.105 п.6.8.7
\hline
\rowcolor{gray!20}
Параметр & Условное обозначение на схеме & Значение/ Диапазон & Единица измерения & Шаг \\ 
%\hline
\endhead

%===m>ЛО внеш АПТ|general+-
\multicolumn{5}{|c|}{ ЛО внеш АПТ } \\ \hline 
\raggedright ЛО внеш АПТ & \centering Введено \\ Выведено & \centering --- & \centering --- & \centering \arraybackslash Введено \\
\hline
%===m
\end{longtable}
\normalsize


%%%%%%%%%%%%%%%%%%%%%%%%%%%%%%%%%%%%%%%%%%%%%%%%%%%%%%%%%%%%%%%%%%%%%%%%% ЛД НН %%%%%%%%%%%%%%%%%%%%%%%%%%%%%%%%%%%%%%%%%%%%%%%%%%%%%%%%%%%%%%%%%%%%%%%%%%%

\par
Уставки ЛД НН представлены в таблице \ref{app_set:ldnn}.

\footnotesize
\begin{longtable}{|>{\centering\arraybackslash}m{5.3cm}|>{\centering\arraybackslash}m{3.3cm}|>{\centering\arraybackslash}m{4.2cm}|>{\centering\arraybackslash}m{1.8cm}|>{\centering\arraybackslash}m{1cm}|}
\caption{Параметры для настройки ЛД НН\hfill\vspace{-0.5\baselineskip}}\label{app_set:ldnn}\\ 
\hline
\rowcolor{gray!20}
Параметр & Условное обозначение на схеме & Значение/ Диапазон & Единица измерения & Шаг \\ 
\hline
\endfirsthead
\caption*{\hspace{3pt}\emph{Продолжение таблицы \ref{app_set:ldnn}\hfill\vspace{-0.5\baselineskip}}} \\ % сделано по ГОСТ 2.105 п.6.8.7
\hline
\rowcolor{gray!20}
Параметр & Условное обозначение на схеме & Значение/ Диапазон & Единица измерения & Шаг \\ 
%\hline
\endhead

%===m>ЛД НН|general+-
\multicolumn{5}{|c|}{ ЛД НН от МТЗ-1 } \\ \hline 
\raggedright Тср МТЗ-1 на В НН & \centering 0 ... 300000 & \centering мс & \centering 5 & \centering \arraybackslash 300000 \\
\hline
\raggedright ЛД НН от МТЗ-1 НН & \centering Введено \\ Выведено & \centering --- & \centering --- & \centering \arraybackslash Введено \\
\hline
\multicolumn{5}{|c|}{ ЛД НН от МТЗ-2 } \\ \hline 
\raggedright Тср МТЗ-2 на В НН & \centering 0 ... 300000 & \centering мс & \centering 5 & \centering \arraybackslash 300000 \\
\hline
\raggedright ЛД НН от МТЗ-2 НН & \centering Введено \\ Выведено & \centering --- & \centering --- & \centering \arraybackslash Введено \\
\hline
%===m
\end{longtable}
\normalsize

%%%%%%%%%%%%%%%%%%%%%%%%%%%%%%%%%%%%%%%%%%%%%%%%%%%%%%%%%%%%%%%%%%%%%%%%% ЛО ОК %%%%%%%%%%%%%%%%%%%%%%%%%%%%%%%%%%%%%%%%%%%%%%%%%%%%%%%%%%%%%%%%%%%%%%%%%%%

\par
Уставки ЛО ОК представлены в таблице \ref{app_set:look}.

\footnotesize
\begin{longtable}{|>{\centering\arraybackslash}m{5.3cm}|>{\centering\arraybackslash}m{3.3cm}|>{\centering\arraybackslash}m{4.2cm}|>{\centering\arraybackslash}m{1.8cm}|>{\centering\arraybackslash}m{1cm}|}
\caption{Параметры для настройки ЛД НН\hfill\vspace{-0.5\baselineskip}}\label{app_set:look}\\ 
\hline
\rowcolor{gray!20}
Параметр & Условное обозначение на схеме & Значение/ Диапазон & Единица измерения & Шаг \\ 
\hline
\endfirsthead
\caption*{\hspace{3pt}\emph{Продолжение таблицы \ref{app_set:look}\hfill\vspace{-0.5\baselineskip}}} \\ % сделано по ГОСТ 2.105 п.6.8.7
\hline
\rowcolor{gray!20}
Параметр & Условное обозначение на схеме & Значение/ Диапазон & Единица измерения & Шаг \\ 
%\hline
\endhead

%===m>ЛД ОК|general+-
\multicolumn{5}{|c|}{ ЛЗ ОК } \\ \hline 
\raggedright Тимп ОК & \centering 0 ... 300000 & \centering мс & \centering 5 & \centering \arraybackslash 1000 \\
\hline
\raggedright Действие на ОК & \centering Введено \\ Выведено & \centering --- & \centering --- & \centering \arraybackslash Введено \\
\hline
%===m
\end{longtable}
\normalsize